% Generated mechanically from the frozen dpa.tex target.
% Do not edit this generated file; edit the bound locale source instead.

% OK, start here.
%

\setcounter{chapter}{22}
\chapter{除幂代数}



\phantomsection
\label{dpa-section-phantom}


\section{引言}
\label{dpa-section-introduction}

\noindent
本章讨论除幂代数及其应用。一本参考书是 \cite{Berthelot}。





\section{除幂}
\label{dpa-section-divided-powers}

\noindent
本节汇集除幂环的一些结果。
我们采用约定 $0! = 1$（空乘积应当等于 $1$）。

\begin{definition}
\label{dpa-definition-divided-powers}
设 $A$ 是环，$I$ 是 $A$ 的理想。若一族映射
$\gamma_n : I \to I$（$n > 0$）对所有 $n \geq 0$、$m > 0$、
$x, y \in I$ 和 $a \in A$ 满足
\begin{enumerate}
\item $\gamma_1(x) = x$，并约定 $\gamma_0(x) = 1$；
\item $\gamma_n(x)\gamma_m(x) = \frac{(n + m)!}{n! m!} \gamma_{n + m}(x)$,
\item $\gamma_n(ax) = a^n \gamma_n(x)$,
\item $\gamma_n(x + y) = \sum_{i = 0, \ldots, n} \gamma_i(x)\gamma_{n - i}(y)$,
\item $\gamma_n(\gamma_m(x)) = \frac{(nm)!}{n! (m!)^n} \gamma_{nm}(x)$.
\end{enumerate}
则称它为 $I$ 上的{\it 除幂结构}。
\end{definition}

\noindent
注意，定义中出现的有理数 $\frac{(n + m)!}{n! m!}$
和 $\frac{(nm)!}{n! (m!)^n}$ 实际上都是整数；
前者是从 $n + m$ 个对象中选取 $n$ 个的方式数，
后者是把 $nm$ 个对象分成 $n$ 个各含 $m$ 个对象的组的方式数。
下面关于该定义的若干说明表明，$\gamma_n(x)$ 在 $I$ 中代替了 $x^n/n!$。

\begin{lemma}
\label{dpa-lemma-silly}
设 $A$ 是环，$I$ 是 $A$ 的理想。
\begin{enumerate}
\item 若 $\gamma$ 是 $I$ 上的除幂结构\footnote{这里以及下文中，
$\gamma$ 是从 $I$ 到 $I$ 的映射序列
$\gamma_1, \gamma_2, \gamma_3, \ldots$ 的简写。}，则对
$n \geq 1$、$x \in I$ 有 $n! \gamma_n(x) = x^n$。
\end{enumerate}
假设 $A$ 作为 $\mathbf{Z}$-模无挠。
\begin{enumerate}
\item[(2)] $I$ 上的除幂结构若存在，则唯一。
\item[(3)] 若 $\gamma_n : I \to I$ 是映射，则
$$
\gamma\text{ 是除幂结构}
\Leftrightarrow
n! \gamma_n(x) = x^n\ \forall x \in I, n \geq 1.
$$
\item[(4)] 理想 $I$ 具有除幂结构，当且仅当存在 $I$ 作为理想的一组生成元
$x_i$，使得对所有 $n \geq 1$ 都有 $x_i^n \in (n!)I$。
\end{enumerate}
\end{lemma}

\begin{proof}
\emph{(1) 的证明。}
若 $\gamma$ 是除幂结构，则条件 (2)（把其中的 $n,m$ 换成
$1,n-1$）给出
$n \gamma_n(x) = \gamma_1(x)\gamma_{n - 1}(x)$。
因此，由归纳法和条件 (1) 得到 $n! \gamma_n(x) = x^n$。

\medskip\noindent
假设 $A$ 作为 $\mathbf{Z}$-模无挠。
\emph{(2) 的证明。} 这由 (1) 立即得到。

\medskip\noindent
\emph{(3) 的证明。}
假设对所有 $x \in I$ 和 $n \geq 1$ 都有
$n! \gamma_n(x) = x^n$。由于
$A \subset A \otimes_{\mathbf{Z}} \mathbf{Q}$，只须在
$A$ 是 $\mathbf{Q}$-代数时证明定义
\ref{dpa-definition-divided-powers} 的公理 (1)--(5)。
此时 $\gamma_n(x) = x^n/n!$，直接验证即可得到 (1)--(5)；
例如，(4) 对应于二项式公式
$$
(x + y)^n = \sum_{i = 0, \ldots, n} \frac{n!}{i!(n - i)!} x^iy^{n - i}
$$
读者不妨完成这些验证，以确认各系数无误。

\medskip\noindent
\emph{(4) 的证明。}
假设 $x_i$ 是 $I$ 作为理想的一组生成元，并且对所有 $n \geq 1$
都有 $x_i^n \in (n!)I$。我们断言，对所有 $x \in I$ 都有
$x^n \in (n!)I$。若该断言成立，则可令
$\gamma_n(x) = x^n/n!$；由 (3)，这给出一个除幂结构。
为证明该断言，注意它对 $x = ax_i$ 成立。因此，该断言对
$I$ 作为阿贝尔群的一组生成元成立。对由这些生成元组成的表达式长度作归纳，
只须证明：若断言对 $x$ 和 $y$ 成立，则它对 $x+y$ 也成立。
这立即由二项式定理得到。
\end{proof}

\begin{example}
\label{dpa-example-ideal-generated-by-p}
设 $p$ 是素数。设 $A$ 是环，并且每个不能被 $p$ 整除的整数 $n$
在 $A$ 中都可逆；也就是说，$A$ 是 $\mathbf{Z}_{(p)}$-代数。
于是 $I = pA$ 具有典范除幂结构。具体地，对给定的
$x = pa \in I$，令
$$
\gamma_n(x) = \frac{p^n}{n!} a^n
$$
读者可以立即验证，对 $n \geq 1$ 有
$p^n/n! \in p\mathbf{Z}_{(p)}$（例如，这可由下述事实推出：
$n!$ 的素因子分解中 $p$ 的指数为
$\left\lfloor n/p \right\rfloor + \left\lfloor n/p^2 \right\rfloor
+ \left\lfloor n/p^3 \right\rfloor + \ldots$）。
因此该定义有意义，并给出一列映射 $\gamma_n : I \to I$。
直接验证即可看出定义 \ref{dpa-definition-divided-powers} 的条件 (1)--(5)
均成立。也可以先注意到该定义对
$A_0 = \mathbf{Z}_{(p)}$ 成立，再由引理 \ref{dpa-lemma-gamma-extends}
得到结论。
\end{example}

\noindent
注意，对 $A$ 的任意理想 $I$ 及其上的任意除幂结构 $\gamma$，
都有 $\gamma_n\left(0\right) = 0$。（这由定义
\ref{dpa-definition-divided-powers} 的公理 (3) 对 $a=0$ 的应用得到。）

\begin{lemma}
\label{dpa-lemma-check-on-generators}
设 $A$ 是环，$I$ 是 $A$ 的理想，并设
$\gamma_n : I \to I$（$n \geq 1$）是一列映射。假设
\begin{enumerate}
\item[(a)] 定义 \ref{dpa-definition-divided-powers} 的 (1)、(3) 和 (4)
对所有 $x,y\in I$ 成立；
\item[(b)] 性质 (2) 和 (5) 对 $I$ 作为理想的某组生成元中的 $x$ 成立。
\end{enumerate}
则 $\gamma$ 是 $I$ 上的除幂结构。
\end{lemma}

\begin{proof}
本证明中的编号 (1)、(2)、(3)、(4)、(5) 指定义
\ref{dpa-definition-divided-powers} 所列的条件。
应用 (3) 可知，若 (2) 和 (5) 对 $x$ 成立，则它们对所有
$a\in A$ 的 $ax$ 也成立。因此，(b) 蕴含 (2) 和 (5)
对 $I$ 作为阿贝尔群的一组生成元成立。
对由这些生成元组成的表达式长度作归纳，只须证明：若
(2) 和 (5) 对 $x,y\in I$ 成立，则它们对 $x+y$ 也成立。

\medskip\noindent
\emph{$x+y$ 的性质 (2)。} 由 (4)，
$$
\gamma_n(x + y)\gamma_m(x + y) =
\sum\nolimits_{i + j = n,\ k + l = m}
\gamma_i(x)\gamma_k(x)\gamma_j(y)\gamma_l(y)
$$
应用 $x$ 和 $y$ 的性质 (2)，上式等于
$$
\sum \frac{(i + k)!}{i!k!}\frac{(j + l)!}{j!l!}
\gamma_{i + k}(x)\gamma_{j + l}(y)
$$
将它与展开式
$$
\gamma_{n + m}(x + y) = \sum \gamma_a(x)\gamma_b(y)
$$
比较可知，只须证明：给定 $a+b=n+m$，有
$$
\sum\nolimits_{i + k = a,\ j + l = b,\ i + j = n,\ k + l = m}
\frac{(i + k)!}{i!k!}\frac{(j + l)!}{j!l!}
=
\frac{(n + m)!}{n!m!}.
$$
不直接证明此式，而是注意到它对多项式环 $\mathbf{Q}[x,y]$
中的理想 $I=(x,y)$ 成立，因为对 $f\in I$，
$\gamma_n(f)=f^n/n!$ 定义了 $I$ 上的除幂结构。
因此，上述有理数恒等式成立。

\medskip\noindent
\emph{$x+y$ 的性质 (5)。}
假设 (1)--(4) 成立，并且 (5) 对 $x$ 和 $y$ 成立。
我们再次把证明归结为一个有理数恒等式。具体地，应用 (4) 可写成
$\gamma_n(\gamma_m(x + y)) = \gamma_n(\sum \gamma_i(x)\gamma_j(y))$.
再次应用 (4)，可把 $\gamma_n(\gamma_m(x+y))$
写成若干项之和，其中每一项都是形如
$\gamma_k(\gamma_i(x)\gamma_j(y))$ 的因子的乘积。
若 $i>0$，则
\begin{align*}
\gamma_k(\gamma_i(x)\gamma_j(y)) & =
\gamma_j(y)^k\gamma_k(\gamma_i(x)) \\
& = \frac{(ki)!}{k!(i!)^k} \gamma_j(y)^k \gamma_{ki}(x) \\
& =
\frac{(ki)!}{k!(i!)^k} \frac{(kj)!}{(j!)^k} \gamma_{ki}(x) \gamma_{kj}(y)
\end{align*}
其中第一个等号使用 (3)，第二个等号使用 $x$ 的 (5)，第三个等号
恰好使用 $k$ 次 (2)。应用 $y$ 的 (5) 可知，该等式在 $i=0$
时同样成立。继续如此应用除 (5) 外的所有公理，可写成
$$
\gamma_n(\gamma_m(x + y)) =
\sum\nolimits_{i + j = nm} c_{ij}\gamma_i(x)\gamma_j(y)
$$
其中 $c_{ij}\in\mathbf{Z}$ 是某些普适常数。该等式在多项式环
$\mathbf{Q}[x,y]$ 中成立这一事实再次蕴含所有系数 $c_{ij}$
均等于 $(nm)!/n!(m!)^n$，正如所需。
\end{proof}

\begin{lemma}
\label{dpa-lemma-two-ideals}
设 $A$ 是环，并有两个理想 $I,J\subset A$。
设 $\gamma$ 是 $I$ 上的除幂结构，$\delta$ 是 $J$ 上的除幂结构。则
\begin{enumerate}
\item $\gamma$ 与 $\delta$ 在 $IJ$ 上一致；
\item 若 $\gamma$ 与 $\delta$ 在 $I\cap J$ 上一致，则它们是
$I+J$ 上唯一除幂结构 $\epsilon$ 的限制。
\end{enumerate}
\end{lemma}

\begin{proof}
设 $x\in I$ 且 $y\in J$。则
$$
\gamma_n(xy) = y^n\gamma_n(x) = n! \delta_n(y) \gamma_n(x) =
\delta_n(y) x^n = \delta_n(xy).
$$
因此，$\gamma$ 与 $\delta$ 在 $IJ$ 的一组（加法）生成元上相同。
由定义 \ref{dpa-definition-divided-powers} 的性质 (4)，它们在整个
$IJ$ 上相同。

\medskip\noindent
假设 $\gamma$ 与 $\delta$ 在 $I\cap J$ 上一致。
给定 $z\in I+J$，将其写成 $z=x+y$，其中 $x\in I$ 且 $y\in J$。
对所有 $n\geq 1$，令
$$
\epsilon_n(z) = \sum \gamma_i(x)\delta_{n - i}(y)
$$
为说明它是良定义的，设 $z=x'+y'$ 是另一种表示，其中
$x'\in I$ 且 $y'\in J$。于是
$w=x-x'=y'-y\in I\cap J$。因此
\begin{align*}
\sum\nolimits_{i + j = n} \gamma_i(x)\delta_j(y)
& =
\sum\nolimits_{i + j = n} \gamma_i(x' + w)\delta_j(y) \\
& =
\sum\nolimits_{i' + l + j = n} \gamma_{i'}(x')\gamma_l(w)\delta_j(y) \\
& =
\sum\nolimits_{i' + l + j = n} \gamma_{i'}(x')\delta_l(w)\delta_j(y) \\
& =
\sum\nolimits_{i' + j' = n} \gamma_{i'}(x')\delta_{j'}(y + w) \\
& =
\sum\nolimits_{i' + j' = n} \gamma_{i'}(x')\delta_{j'}(y')
\end{align*}
正如所需。因此，我们对所有 $n\geq 1$ 定义了映射
$\epsilon_n:I+J\to I+J$；容易看出
$\epsilon_n\mid_I=\gamma_n$ 且 $\epsilon_n\mid_J=\delta_n$。
下面验证映射族 $\epsilon_n$ 满足定义
\ref{dpa-definition-divided-powers} 的条件 (1)--(5)。
性质 (1) 和 (3) 是显然的。为验证 (4)，设
$z=x+y$、$z'=x'+y'$，其中 $x,x'\in I$ 且 $y,y'\in J$，计算
\begin{align*}
\epsilon_n(z + z') & =
\sum\nolimits_{a + b = n} \gamma_a(x + x')\delta_b(y + y') \\
& =
\sum\nolimits_{i + i' + j + j' = n}
\gamma_i(x) \gamma_{i'}(x')\delta_j(y)\delta_{j'}(y') \\
& =
\sum\nolimits_{k = 0, \ldots, n}
\sum\nolimits_{i+j=k} \gamma_i(x)\delta_j(y)
\sum\nolimits_{i'+j'=n-k} \gamma_{i'}(x')\delta_{j'}(y') \\
& =
\sum\nolimits_{k = 0, \ldots, n}\epsilon_k(z)\epsilon_{n-k}(z')
\end{align*}
正如所需。现在只须对 $I$ 或 $J$ 中的元素证明 (2) 和 (5)，
见引理 \ref{dpa-lemma-check-on-generators}。
由于 $\gamma$ 和 $\delta$ 都是除幂结构，这一点是显然的。

\medskip\noindent
由此证明了 $I+J$ 上存在除幂结构 $\epsilon$，其在 $I$ 和 $J$
上的限制分别为 $\gamma$ 和 $\delta$；其唯一性也很清楚。
\end{proof}

\begin{lemma}
\label{dpa-lemma-nil}
设 $p$ 是素数，$A$ 是环，$I\subset A$ 是理想，并设
$\gamma$ 是 $I$ 上的除幂结构。假设 $p$ 在 $A/I$ 中幂零。
则 $I$ 局部幂零，当且仅当 $p$ 在 $A$ 中幂零。
\end{lemma}

\begin{proof}
若在 $A$ 中 $p^N=0$，则对 $x\in I$ 有
$x^{pN}=(pN)!\gamma_{pN}(x)=0$，因为 $(pN)!$ 可被 $p^N$ 整除。
反之，假设 $I$ 局部幂零。我们还假设 $p$ 在 $A/I$ 中幂零，
故对某个 $r$ 有 $p^r\in I$；于是 $p^r$ 幂零，从而 $p$ 幂零。
\end{proof}








\section{除幂环}
\label{dpa-section-divided-power-rings}

\noindent
除幂环构成一个范畴。定义如下。

\begin{definition}
\label{dpa-definition-divided-power-ring}
{\it 除幂环}是三元组 $(A,I,\gamma)$，其中 $A$ 是环，
$I\subset A$ 是理想，而
$\gamma=(\gamma_n)_{n\geq 1}$ 是 $I$ 上的除幂结构。
{\it 除幂环同态}
$\varphi:(A,I,\gamma)\to(B,J,\delta)$ 是环同态
$\varphi:A\to B$，满足 $\varphi(I)\subset J$，并且对所有
$x\in I$ 和 $n\geq 1$ 都有
$\delta_n(\varphi(x))=\varphi(\gamma_n(x))$。
\end{definition}

\noindent
有时我们说“设 $(B,J,\delta)$ 是 $(A,I,\gamma)$ 上的除幂代数”，
意思是 $(B,J,\delta)$ 是一个除幂环，并配备了一个除幂环同态
$(A,I,\gamma)\to(B,J,\delta)$。

\begin{lemma}
\label{dpa-lemma-limits}
除幂环范畴具有所有极限，并且这些极限与环范畴中的极限一致。
\end{lemma}

\begin{proof}
空极限是零环（这有些奇怪，但我们需要如此）。
一族除幂环 $(A_t,I_t,\gamma_t)$（$t\in T$）的积为
$(\prod A_t,\prod I_t,\gamma)$，其中
$\gamma_n((x_t))=(\gamma_{t,n}(x_t))$。
两个同态
$\alpha,\beta:(A,I,\gamma)\to(B,J,\delta)$ 的等化子就是
$C=\{a\in A\mid\alpha(a)=\beta(a)\}$，配备理想 $C\cap I$
及其诱导的除幂。由此所有极限都存在，见
《范畴》引理 \ref{categories-lemma-limits-products-equalizers}。
\end{proof}

\noindent
下述引理以除幂代数为例说明一种很一般的范畴论现象。

\begin{lemma}
\label{dpa-lemma-a-version-of-brown}
设 $\mathcal{C}$ 是除幂环范畴，并设
$F:\mathcal{C}\to\textit{Sets}$ 是函子。假设
\begin{enumerate}
\item 存在基数 $\kappa$，使得对每个 $f\in F(A,I,\gamma)$，
都存在 $\mathcal{C}$ 中的态射
$(A',I',\gamma')\to(A,I,\gamma)$，其中 $|A'|\leq\kappa$，
并且 $f$ 是某个 $f'\in F(A',I',\gamma')$ 的像；
\item $F$ 与极限交换。
\end{enumerate}
则 $F$ 可表；也就是说，存在 $\mathcal{C}$ 的对象 $(B,J,\delta)$，
使得
$$
F(A, I, \gamma) = \Hom_\mathcal{C}((B, J, \delta), (A, I, \gamma))
$$
关于 $(A,I,\gamma)$ 函子性地成立。
\end{lemma}

\begin{proof}
这是《范畴》引理 \ref{categories-lemma-a-version-of-brown}
的一个特殊情形。
\end{proof}

\begin{lemma}
\label{dpa-lemma-colimits}
除幂环范畴具有所有余极限。
\end{lemma}

\begin{proof}
空余极限是带除幂理想 $(0)$ 的 $\mathbf{Z}$。
下面讨论一般的余极限。设 $\mathcal{C}$ 是范畴，并设
$c\mapsto(A_c,I_c,\gamma_c)$ 是一个图。考虑函子
$$
F(B, J, \delta) = \lim_{c \in \mathcal{C}}
Hom((A_c, I_c, \gamma_c), (B, J, \delta))
$$
注意，任意
$f=(f_c)_{c\in C}\in F(B,J,\delta)$ 都具有如下性质：
所有像 $f_c(A_c)$ 生成 $B$ 的一个基数受 $\kappa$ 控制的子环
$B'$，而所有像 $f_c(I_c)$ 在 $B'$ 中生成一个除幂子理想 $J'$。
于是 $f$ 分解为 $F(B')$ 中的某个 $f'$，再复合包含
$B'\to B$。此外，$F$ 与极限交换。因此可以应用引理
\ref{dpa-lemma-a-version-of-brown}，得到 $F$ 可表，结论成立。
\end{proof}

\begin{remark}
\label{dpa-remark-pushout}
设
$$
\xymatrix{
(B, J, \delta) \ar[r] & (B'', J'', \delta'') \\
(A, I, \gamma) \ar[r] \ar[u] & (B', J', \delta') \ar[u]
}
$$
是除幂环范畴中的一个推出（推出由引理 \ref{dpa-lemma-colimits} 存在）。则
\begin{enumerate}
\item $B''/J'' = B/J \otimes_{A/I} B'/J'$,
\item 映射 $\varphi : B \otimes_A B' \to B''$ 是满射；
\item 由 $\varphi$ 诱导的映射
$J \otimes_A B' + B \otimes_A J' \to J''$ 是满射。
\end{enumerate}
为证明 (1)，考虑映射
$(B'', J'', \delta'') \to (C, (0), \emptyset)$，并应用推出的定义。
为证明 (2)，考虑像 $\tilde B = \varphi(B \otimes_A B')$ 以及理想
$\tilde J = \varphi(J \otimes_A B' + B \otimes_A J')$。
由于 $\delta''$ 与 $\delta$、$\delta'$ 相容，显然对所有 $n \geq 1$ 都有
$\delta''_n(\tilde J) \subset \tilde J$。
因此，由推出的泛性质必有 $\tilde B = B''$ 且 $\tilde J = J''$。
\end{remark}

\begin{remark}
\label{dpa-remark-forgetful}
忘却函子 $(A, I, \gamma) \mapsto A$ 不与余极限交换。例如，设
$$
\xymatrix{
(B, J, \delta) \ar[r] & (B'', J'', \delta'') \\
(A, I, \gamma) \ar[r] \ar[u] & (B', J', \delta') \ar[u]
}
$$
是评注 \ref{dpa-remark-pushout} 中所讨论的除幂环范畴中的推出。
那么一般而言，映射 $B \otimes_A B' \to B''$ 不是同构。
一个具体例子如下：
$(A, I, \gamma) = (\mathbf{Z}, (0), \emptyset)$,
$(B, J, \delta) = (\mathbf{Z}/4\mathbf{Z}, 2\mathbf{Z}/4\mathbf{Z}, \delta)$,
以及
$(B', J', \delta') =
(\mathbf{Z}/4\mathbf{Z}, 2\mathbf{Z}/4\mathbf{Z}, \delta')$
其中 $\delta_2(2) = 2$，而 $\delta'_2(2) = 0$。
更准确地说，利用引理 \ref{dpa-lemma-need-only-gamma-p}，分别令
$\delta$、$\delta'$ 为 $J$、$J'$ 上唯一的除幂结构，使得
$\delta_2 : J \to J$、$\delta'_2 : J' \to J'$ 分别是映射
$0 \mapsto 0, 2 \mapsto 2$ 和 $0 \mapsto 0, 2 \mapsto 0$。
于是 $(B'', J'', \delta'') = (\mathbf{F}_2, (0), \emptyset)$，
这与张量积并不一致。
\end{remark}


\section{除幂结构的延拓}
\label{dpa-section-extend}

\noindent
定义如下。

\begin{definition}
\label{dpa-definition-extends}
给定除幂环 $(A, I, \gamma)$ 和环同态 $A \to B$，若 $IB$ 上存在
除幂结构 $\bar \gamma$，使得
$(A, I, \gamma) \to (B, IB, \bar\gamma)$ 是除幂环同态，
则称 $\gamma$ 可{\it 延拓}到 $B$。
\end{definition}

\begin{lemma}
\label{dpa-lemma-gamma-extends}
设 $(A, I, \gamma)$ 是除幂环，$A \to B$ 是环同态。
若 $\gamma$ 可延拓到 $B$，则其延拓唯一。
假设下列条件中至少有一个成立：
\begin{enumerate}
\item $IB = 0$,
\item $I$ 是主理想；
\item $A \to B$ 是平坦的。
\end{enumerate}
则 $\gamma$ 可延拓到 $B$。
\end{lemma}

\begin{proof}
$IB$ 的任意元素都可写成有限和
$\sum\nolimits_{i=1}^t b_ix_i$，其中 $b_i \in B$、$x_i \in I$。
若 $\gamma$ 延拓为 $IB$ 上的 $\bar\gamma$，则
$\bar\gamma_n(x_i) = \gamma_n(x_i)$。
因此，定义 \ref{dpa-definition-divided-powers} 中的条件 (3) 和 (4) 蕴含
$$
\bar\gamma_n(\sum\nolimits_{i=1}^t b_ix_i) =
\sum\nolimits_{n_1 + \ldots + n_t = n}
\prod\nolimits_{i = 1}^t b_i^{n_i}\gamma_{n_i}(x_i)
$$
由此可见，$\bar\gamma$ 若存在便是唯一的。

\medskip\noindent
若 $IB = 0$，令 $\bar\gamma_n(0) = 0$ 即可。若 $I = (x)$，
则定义 $\bar\gamma_n(bx) = b^n\gamma_n(x)$。这是良定义的：
若 $b'x = bx$，即 $(b - b')x = 0$，则
\begin{align*}
b^n\gamma_n(x) - (b')^n\gamma_n(x)
& =
(b^n - (b')^n)\gamma_n(x) \\
& =
(b^{n - 1} + \ldots + (b')^{n - 1})(b - b')\gamma_n(x) = 0
\end{align*}
因为 $\gamma_n(x)$ 可被 $x$ 整除（这是由于
$\gamma_n(I) \subset I$），从而被 $b-b'$ 零化。
下面证明定义 \ref{dpa-definition-divided-powers} 的条件 (1)--(5)。
(1)、(2)、(3)、(5) 由构造显然成立。
对于 (4)，设 $y,z\in IB$，写成 $y=bx$、$z=cx$。
则 $y+z=(b+c)x$，因而
\begin{align*}
\bar\gamma_n(y + z)
& =
(b + c)^n\gamma_n(x) \\
& =
\sum \frac{n!}{i!(n - i)!}b^ic^{n -i}\gamma_n(x) \\
& =
\sum b^ic^{n - i}\gamma_i(x)\gamma_{n - i}(x) \\
& =
\sum \bar\gamma_i(y)\bar\gamma_{n -i}(z)
\end{align*}
正是所需结论。

\medskip\noindent
假设 $A \to B$ 是平坦的。设 $b_1, \ldots, b_r \in B$ 且
$x_1, \ldots, x_r \in I$。若 $\bar\gamma_n$ 存在，则
$$
\bar\gamma_n(\sum b_ix_i) =
\sum b_1^{e_1} \ldots b_r^{e_r} \gamma_{e_1}(x_1) \ldots \gamma_{e_r}(x_r)
$$
其中求和取遍 $e_1 + \ldots + e_r = n$。
再设有 $c_1, \ldots, c_s \in B$ 和 $a_{ij} \in A$，使得
$b_i = \sum a_{ij}c_j$。令 $y_j = \sum a_{ij}x_i$，我们断言
$$
\sum b_1^{e_1} \ldots b_r^{e_r} \gamma_{e_1}(x_1) \ldots \gamma_{e_r}(x_r) =
\sum c_1^{d_1} \ldots c_s^{d_s} \gamma_{d_1}(y_1) \ldots \gamma_{d_s}(y_s)
$$
在 $B$ 中成立，其中右端求和取遍 $d_1 + \ldots + d_s = n$。
事实上，利用除幂结构的公理，可将两端都展开成若干项之和；
其系数属于 $\mathbf{Z}[a_{ij}]$，各项形如
$c_1^{d_1} \ldots c_s^{d_s}\gamma_{e_1}(x_1) \ldots \gamma_{e_r}(x_r)$。
为说明系数相同，注意当
$\gamma$ 是 $(x_1, \ldots, x_r)$ 上唯一的除幂结构时，
该结论在 $\mathbf{Q}[x_1, \ldots, x_r, c_1, \ldots, c_s, a_{ij}]$
中成立。由 Lazard 定理（《代数》定理 \ref{algebra-theorem-lazard}），
可将 $B$ 写成有限自由 $A$-模的有向余极限。
特别地，若 $z\in IB$ 既写成 $z=\sum x_ib_i$，又写成
$z=\sum x'_{i'}b'_{i'}$，则可找到 $c_1,\ldots,c_s\in B$ 以及
$a_{ij},a'_{i'j}\in A$，使得 $b_i=\sum a_{ij}c_j$、
$b'_{i'}=\sum a'_{i'j}c_j$，并且
$y_j = \sum x_ia_{ij} = \sum x'_{i'}a'_{i'j}$\footnote{也可以不借助
《代数》定理 \ref{algebra-theorem-lazard} 来证明这一点。事实上，若
$z = \sum x_ib_i$ 且 $z = \sum x'_{i'}b'_{i'}$，则
$\sum x_ib_i - \sum x'_{i'}b'_{i'} = 0$ 是 $A$-模 $B$ 中的一个关系。
因此，《代数》引理 \ref{algebra-lemma-flat-eq}（以 $x_i$ 和
$x'_{i'}$ 代替其中的 $f_i$，并以 $b_i$ 和 $b'_{i'}$ 代替其中的
$x_i$）给出所需的 $c_1, \ldots, c_s \in B$ 以及
$a_{ij}, a'_{i'j} \in A$。}。
因此，上述过程给出 $IB$ 上一个良定义的映射 $\bar\gamma_n$。
由构造，$\bar\gamma$ 满足条件 (1)、(3)、(4)。此外，对 $x\in I$ 有
$\bar\gamma_n(x)=\gamma_n(x)$。故由引理
\ref{dpa-lemma-check-on-generators}，$\bar\gamma$ 是 $IB$ 上的除幂结构。
\end{proof}

\begin{lemma}
\label{dpa-lemma-kernel}
设 $(A, I, \gamma)$ 是除幂环。
\begin{enumerate}
\item 若 $\varphi : (A, I, \gamma) \to (B, J, \delta)$ 是除幂环同态，
则对所有 $n\geq 1$，$\Ker(\varphi) \cap I$ 在 $\gamma_n$ 下不变。
\item 设 $\mathfrak a \subset A$ 是理想，并令
$I' = I \cap \mathfrak a$。下列条件等价：
\begin{enumerate}
\item 对所有 $n > 0$，$I'$ 在 $\gamma_n$ 下不变；
\item $\gamma$ 可延拓到 $A/\mathfrak a$；
\item 存在 $I'$ 作为理想的一组生成元 $x_i$，使得对所有 $n>0$，
$\gamma_n(x_i) \in I'$。
\end{enumerate}
\end{enumerate}
\end{lemma}

\begin{proof}
证明 (1)。这是显然的。假设 (2)(a)。对 $x\in I$，定义
$\bar\gamma_n(x \bmod I') = \gamma_n(x) \bmod I'$。
这是良定义的，因为由定义 \ref{dpa-definition-divided-powers} (4) 以及
假设中的 $\gamma_j(y)\in I'$，对 $y\in I'$ 有
$\gamma_n(x+y)=\gamma_n(x)\bmod I'$。由于 $\gamma$ 是除幂结构，
$\bar\gamma$ 显然也是除幂结构。因此 (2)(b) 成立。
另一方面，由 (1)，(2)(b) 蕴含 (2)(a)。显然 (2)(a) 蕴含 (2)(c)。
现在假设 (2)(c)。注意当 $x=ax_i$ 时，
$\gamma_n(x)=a^n\gamma_n(x_i)\in I'$。
因此，对 $I'$ 作为阿贝尔群的一组生成元，均有
$\gamma_n(x)\in I'$。对由这些生成元组成的表达式长度作归纳，
只须证明：若 $\forall n : \gamma_n(x), \gamma_n(y) \in I'$，则
$\forall n : \gamma_n(x+y)\in I'$。这立即由除幂结构的第四条公理得到。
\end{proof}

\begin{lemma}
\label{dpa-lemma-sub-dp-ideal}
设 $(A, I, \gamma)$ 是除幂环，$E\subset I$ 是子集。
则在 $\gamma$ 下不变且包含所有 $f\in E$ 的最小理想 $J\subset I$，
就是由 $\gamma_n(f)$（$n\geq 1$、$f\in E$）生成的理想 $J$。
\end{lemma}

\begin{proof}
立即由引理 \ref{dpa-lemma-kernel} 得到。
\end{proof}

\begin{lemma}
\label{dpa-lemma-extend-to-completion}
设 $(A, I, \gamma)$ 是除幂环，$p$ 是素数。若 $p$ 在 $A/I$ 中幂零，则
\begin{enumerate}
\item $p$-进完备化 $A^\wedge = \lim_e A/p^eA$ 满射到 $A/I$；
\item 该映射的核是 $I$ 的 $p$-进完备化 $I^\wedge$；
\item 每个 $\gamma_n$ 关于 $p$-进拓扑连续，并延拓为
$\gamma_n^\wedge : I^\wedge \to I^\wedge$，从而在 $I^\wedge$ 上定义
一个除幂结构。
\end{enumerate}
若进一步假设 $A$ 是 $\mathbf{Z}_{(p)}$-代数，则
\begin{enumerate}
\item[(4)] 当 $e$ 足够大时，理想 $p^eA \subset I$ 在除幂结构
$\gamma$ 下不变，并且
$$
(A^\wedge, I^\wedge, \gamma^\wedge) = \lim_e (A/p^eA, I/p^eA, \bar\gamma)
$$
在除幂环范畴中成立。
\end{enumerate}
\end{lemma}

\begin{proof}
取整数 $t\geq 1$，使得 $p^tA/I=0$，即 $p^tA\subset I$。
映射 $A^\wedge\to A/I$ 是复合
$A^\wedge\to A/p^tA\to A/I$，它是满射（例如由
《代数》引理 \ref{algebra-lemma-completion-generalities}）。
对 $e\geq t$，由
$p^eI \subset p^eA \cap I \subset p^{e - t}I$ 可见，
复合 $A^\wedge\to A/I$ 的核是 $I$ 的 $p$-进完备化。
映射 $\gamma_n$ 连续，因为
$$
\gamma_n(x + p^ey) =
\sum\nolimits_{i + j = n} p^{je}\gamma_i(x)\gamma_j(y) =
\gamma_n(x) \bmod p^eI
$$
由除幂结构的公理即得。显然，映射 $\gamma_n^\wedge$ 从
$\gamma_n$ 继承除幂结构的各条公理。最后证明末一断言，取 $e>t$。
此时 $p^eA\subset I$ 且 $\gamma_1(p^eA)\subset p^eA$；对 $n>1$，有
$$
\gamma_n(p^ea) = p^n \gamma_n(p^{e - 1}a) = \frac{p^n}{n!} p^{n(e - 1)}a^n
\in p^e A
$$
因为 $p^n/n! \in \mathbf{Z}_{(p)}$，并且 $n\geq 2$、$e\geq 2$，
故 $n(e-1)\geq e$。
这证明 $\gamma$ 可延拓到 $A/p^eA$，参见引理 \ref{dpa-lemma-kernel}。
关于极限的断言由引理 \ref{dpa-lemma-limits} 的证明中对极限的构造显然成立。
\end{proof}




\section{除幂多项式代数}
\label{dpa-section-divided-power-polynomial-ring}

\noindent
一个非常有用的例子是{\it 除幂多项式代数}。
设 $A$ 是环，$t\geq 1$。用
$A\langle x_1, \ldots, x_t \rangle$ 表示如下 $A$-代数：
作为 $A$-模，令
$$
A\langle x_1, \ldots, x_t \rangle =
\bigoplus\nolimits_{n_1, \ldots, n_t \geq 0} A x_1^{[n_1]} \ldots x_t^{[n_t]}
$$
其乘法由下式给出：
$$
x_i^{[n]}x_i^{[m]} = \frac{(n + m)!}{n!m!}x_i^{[n + m]}.
$$
另令 $x_i=x_i^{[1]}$。注意
$1=x_1^{[0]}\ldots x_t^{[0]}$。类似的构造给出具有无限多个变量的
除幂多项式代数。存在典范 $A$-代数同态
$A\langle x_1, \ldots, x_t \rangle \to A$，它在 $n>0$ 时将
$x_i^{[n]}$ 映为零。该映射的核记为
$A\langle x_1, \ldots, x_t \rangle_{+}$。

\begin{lemma}
\label{dpa-lemma-divided-power-polynomial-algebra}
设 $(A,I,\gamma)$ 是除幂环。下列理想上存在唯一的除幂结构 $\delta$：
$$
J = IA\langle x_1, \ldots, x_t \rangle + A\langle x_1, \ldots, x_t \rangle_{+}
$$
使得
\begin{enumerate}
\item $\delta_n(x_i) = x_i^{[n]}$；
\item $(A, I, \gamma) \to (A\langle x_1, \ldots, x_t \rangle, J, \delta)$
是除幂环同态。
\end{enumerate}
此外，$(A\langle x_1, \ldots, x_t \rangle, J, \delta)$ 具有如下
泛性质：一个除幂环同态
$\varphi : (A\langle x_1, \ldots, x_t \rangle, J, \delta) \to
(C, K, \epsilon)$
与一个除幂环同态 $A\to C$ 以及元素
$k_1, \ldots, k_t \in K$ 等价。
\end{lemma}

\begin{proof}
我们在单变量除幂多项式代数的情形证明该引理。一般情形可完全类似地证明；
也可以注意到，$A\langle x_1, \ldots, x_t\rangle$ 同构于逐个添加
除幂变量 $x_1,\ldots,x_t$ 所得到的环。

\medskip\noindent
设 $A\langle x \rangle_{+}$ 是由
$x, x^{[2]}, x^{[3]}, \ldots$ 生成的理想。
注意 $J = IA\langle x \rangle + A\langle x \rangle_{+}$，并且
$$
IA\langle x \rangle \cap A\langle x \rangle_{+} =
IA\langle x \rangle \cdot A\langle x \rangle_{+}
$$
故由引理 \ref{dpa-lemma-two-ideals}，只须证明理想
$IA\langle x \rangle$ 和 $A\langle x \rangle_{+}$ 上存在除幂结构。
由于 $A\to A\langle x\rangle$ 平坦，前者的存在性由引理
\ref{dpa-lemma-gamma-extends} 得到。对于后者，注意若 $A$ 无挠，
则可应用引理 \ref{dpa-lemma-silly} (4) 得到 $\delta$ 的存在性。
事实上，选取元素 $x^{[m]}$ 作为生成元，可见
$(x^{[m]})^n = \frac{(nm)!}{(m!)^n} x^{[nm]}$
且 $n!$ 整除整数 $\frac{(nm)!}{(m!)^n}$。
一般而言，将 $A$ 写成某个无挠环 $R$ 的商 $R/\mathfrak a$
（例如可取 $\mathbf{Z}$ 上的多项式环）。映射
$R\langle x \rangle \to A\langle x \rangle$ 的核是
$\bigoplus \mathfrak a x^{[m]}$。应用引理 \ref{dpa-lemma-kernel}
的判据 (2)(c)，可见 $R\langle x \rangle_{+}$ 上的除幂结构
按所需方式延拓到 $A\langle x \rangle$。

\medskip\noindent
证明泛性质。给定除幂环同态 $\varphi:A\to C$ 以及
$k_1,\ldots,k_t\in K$，考虑
$$
A\langle x_1, \ldots, x_t \rangle \to C,\quad
x_1^{[n_1]} \ldots x_t^{[n_t]} \longmapsto
\epsilon_{n_1}(k_1) \ldots \epsilon_{n_t}(k_t)
$$
并在系数上使用 $\varphi$。只须检验这是 $A$-代数同态（细节从略）。
反向构造是显然的。
\end{proof}

\begin{remark}
\label{dpa-remark-divided-power-polynomial-algebra}
设 $(A,I,\gamma)$ 是除幂环。引理
\ref{dpa-lemma-divided-power-polynomial-algebra} 有一个无限多个变量的版本。
首先注意，若 $s<t$，则存在典范映射
$$
A\langle x_1, \ldots, x_s \rangle \to A\langle x_1, \ldots, x_t\rangle
$$
因此，对任意集合 $W$，令
$$
A\langle x_w: w \in W\rangle =
\colim_{E \subset W} A\langle x_e:e \in E\rangle
$$
（余极限取遍 $W$ 的有限子集 $E$），转移映射如上。
由余极限的定义可见，$A\langle x_w:w\in W\rangle$ 的泛映射性质
与引理 \ref{dpa-lemma-divided-power-polynomial-algebra} 中所述的映射性质
完全类似。
\end{remark}

\noindent
下述引理可见于 \cite{BO}。

\begin{lemma}
\label{dpa-lemma-need-only-gamma-p}
设 $p$ 是素数。设 $A$ 是环，并且每个不能被 $p$ 整除的整数 $n$
在 $A$ 中都可逆；也就是说，$A$ 是 $\mathbf{Z}_{(p)}$-代数。
设 $I\subset A$ 是理想。$I$ 上的两个除幂结构
$\gamma,\gamma'$ 相等，当且仅当 $\gamma_p=\gamma'_p$。
此外，给定映射 $\delta:I\to I$，满足
\begin{enumerate}
\item 对所有 $x\in I$，$p!\delta(x)=x^p$\footnote{此条件可由另外两个
条件推出，参见评注 \ref{dpa-remark-ryo-suzuki}。}；
\item 对所有 $a\in A$、$x\in I$，$\delta(ax)=a^p\delta(x)$；
\item
$\delta(x + y) =
\delta(x) +
\sum\nolimits_{i + j = p, i,j \geq 1} \frac{1}{i!j!} x^i y^j +
\delta(y)$ 对所有 $x,y\in I$ 成立，
\end{enumerate}
则 $I$ 上存在唯一的除幂结构 $\gamma$，使得 $\gamma_p=\delta$。
\end{lemma}

\begin{proof}
若 $n$ 不能被 $p$ 整除，则
$\gamma_n(x)=c x\gamma_{n-1}(x)$，其中 $c$ 是
$\mathbf{Z}_{(p)}$ 中的单位。此外，
$$
\gamma_{pm}(x) = c \gamma_m(\gamma_p(x))
$$
其中 $c$ 是 $\mathbf{Z}_{(p)}$ 中的单位。因此第一条断言显然。
对于第二条断言，可反向利用这些等式定义所有 $\gamma_n$。
更准确地说，若
$n=a_0+a_1p+\ldots+a_ep^e$，其中 $a_i\in\{0,\ldots,p-1\}$，则令
$$
\gamma_n(x) = c_n x^{a_0} \delta(x)^{a_1} \ldots \delta^e(x)^{a_e}
$$
其中 $c_n\in\mathbf{Z}_{(p)}$ 定义为
$$
c_n =
{(p!)^{a_1 + a_2(1 + p) + \ldots + a_e(1 + \ldots + p^{e - 1})}}/{n!}.
$$
现在须证明除幂结构的公理 (1)--(5)，参见定义
\ref{dpa-definition-divided-powers}。公理 (1) 和 (3) 是立即可得的；
公理 (2) 和 (5) 可直接计算验证，细节从略。设 $x,y\in I$。
我们断言存在环同态
$$
\varphi : \mathbf{Z}_{(p)}\langle u, v \rangle \longrightarrow A
$$
将 $u^{[n]}$ 映为 $\gamma_n(x)$，将 $v^{[n]}$ 映为 $\gamma_n(y)$。
由 $\mathbf{Z}_{(p)}\langle u,v\rangle$ 的构造，这意味着须检验
$$
\gamma_n(x)\gamma_m(x) = \frac{(n + m)!}{n!m!} \gamma_{n + m}(x)
$$
在 $A$ 中成立，并对 $y$ 也成立。由于 $\gamma$ 满足 (2)，这确实成立。
令 $\epsilon$ 表示 $\mathbf{Z}_{(p)}\langle u,v\rangle$ 的理想
$\mathbf{Z}_{(p)}\langle u,v\rangle_{+}$ 上的除幂结构。
接着断言：对 $f\in\mathbf{Z}_{(p)}\langle u,v\rangle_{+}$ 和所有 $n$，
$\varphi(\epsilon_n(f))=\gamma_n(\varphi(f))$。
当 $n=0,1,\ldots,p-1$ 时这是显然的。当 $n=p$ 时，只须对理想
$\mathbf{Z}_{(p)}\langle u,v\rangle_{+}$ 的一组生成元证明，
因为 $\epsilon_p$ 和 $\gamma_p=\delta$ 都满足该引理的条件 (1)、(3)。
因此，只须证明
$\gamma_p(\gamma_n(x))=\frac{(pn)!}{p!(n!)^p}\gamma_{pn}(x)$，
并对 $y$ 也证明同样的等式；这由 $\gamma$ 满足 (5) 得到。
现在，若任意整数 $n=a_0+a_1p+\ldots+a_ep^e$ 按上述方式写成
$p$-进展开，则
$$
\epsilon_n(f) =
c_n f^{a_0} \gamma_p(f)^{a_1} \ldots \gamma_p^e(f)^{a_e}
$$
因为 $\epsilon$ 是除幂结构。因此，对所有 $n$ 都有
$\varphi(\epsilon_n(f))=\gamma_n(\varphi(f))$。
将其应用于 $f=u+v$，由 $\epsilon$ 是除幂结构可知
$\gamma$ 满足公理 (4)。
\end{proof}

\begin{remark}
\label{dpa-remark-ryo-suzuki}
设 $A\supset I$ 如引理 \ref{dpa-lemma-need-only-gamma-p} 中所述，
并设映射 $\delta:I\to I$ 满足该引理的 (2)、(3)。则 (1) 也成立。
首先，对 $n$ 作归纳，证明当 $n\geq1$ 时
$\delta(nx)=n\delta(x)+\frac{n^p-n}{p!}x^p$。$n=1$ 的情形成立。
假设它对某个 $n$ 成立，则
$\delta((n + 1)x) = \delta(nx) + \delta(x) +
(\sum_{i = 1}^{p - 1} \frac{n^i}{i!(p - i)!})x^p
= (n + 1)\delta(x) + \left( \frac{n^p - n}{p!} + \frac{(n + 1)^p}{p!}
- \frac{n^p + 1}{p!} \right)x^p =
(n + 1)\delta(x) + \frac{(n + 1)^p - (n + 1)}{p!}x^p$,
故它对 $n+1$ 也成立。另一方面，$\delta(nx)=n^p\delta(x)$，所以
$(n^p - n)\delta(x) = \frac{n^p - n}{p!}x^p$. 
但对任意 $p$，都存在 $n$ 使得 $p^2 \not \mid n^p-n$。
事实上，
$(\mathbf{Z}/p^2\mathbf{Z})^\times \cong \mathbf{Z}/p(p-1)\mathbf{Z}$.
因此存在 $n$，使得 $n^{p-1}\not\equiv1\mod p^2$。
于是得到 $p!\delta(x)=x^p$。
\end{remark}









\section{Tate 分辨}
\label{dpa-section-tate}

\noindent
本节简要讨论 \cite{Tate-homology} 和 \cite{AH} 中构造的分辨；
这些分辨将除幂结构与微分分次代数结合起来。
本节对微分分次代数使用{\it 同调记号}。
我们的微分分次代数位于非负同调次数。因此，微分分次代数
$(A,\text{d})$ 由如下链复形给出：
$$
\ldots \to A_2 \to A_1 \to A_0 \to 0 \to \ldots
$$
并配备一个乘法。

\medskip\noindent
设 $R$ 是环（照例假设交换）。本节常考虑负次数分量为零的分次
$R$-代数 $A=\bigoplus_{d\geq0}A_d$。令
$A_+=\bigoplus_{d>0}A_d$，并记
$A_{even}=\bigoplus_{d\geq0}A_{2d}$、
$A_{odd}=\bigoplus_{d\geq0}A_{2d+1}$。
回顾：若对齐次元素 $x,y$ 有
$xy=(-1)^{\deg(x)\deg(y)}yx$，则称 $A$ 分次交换。
若进一步对奇次数齐次元素 $x$ 有 $x^2=0$，则称 $A$ 严格分次交换。
最后，为理解下面的定义，请记住：若 $A$ 是 $\mathbf{Q}$-代数，
则 $\gamma_n(x)=x^n/n!$。

\begin{definition}
\label{dpa-definition-divided-powers-graded}
设 $R$ 是环，$A=\bigoplus_{d\geq0}A_d$ 是严格分次交换的分次
$R$-代数。若对所有 $n>0$ 定义的一族映射
$\gamma_n:A_{even,+}\to A_{even,+}$ 满足
\begin{enumerate}
\item 若 $x\in A_{2d}$，则 $\gamma_n(x)\in A_{2nd}$；
\item 对任意 $x$，$\gamma_1(x)=x$，并令 $\gamma_0(x)=1$；
\item $\gamma_n(x)\gamma_m(x) = \frac{(n + m)!}{n! m!} \gamma_{n + m}(x)$,
\item 对所有 $x\in A_{even}$、$y\in A_{even,+}$，
$\gamma_n(xy)=x^n\gamma_n(y)$；
\item 若 $x,y\in A_{odd}$ 齐次且 $n>1$，则 $\gamma_n(xy)=0$；
\item 若 $x,y\in A_{even,+}$，则
$\gamma_n(x+y)=\sum_{i=0,\ldots,n}\gamma_i(x)\gamma_{n-i}(y)$；
\item $\gamma_n(\gamma_m(x)) =
\frac{(nm)!}{n! (m!)^n} \gamma_{nm}(x)$ 对 $x\in A_{even,+}$ 成立。
\end{enumerate}
则称其为 $A$ 上的{\it 除幂结构}。
\end{definition}

\noindent
注意，条件 (2)、(3)、(4)、(6)、(7) 蕴含 $\gamma$ 是交换环
$A_{even}$ 的理想 $A_{even,+}$ 上的“通常”除幂结构，参见第
\ref{dpa-section-divided-powers}、
\ref{dpa-section-divided-power-rings}、
\ref{dpa-section-extend} 和
\ref{dpa-section-divided-power-polynomial-ring} 节。
特别地，对所有 $x\in A_{even,+}$ 有 $n!\gamma_n(x)=x^n$。
条件 (1) 表明 $\gamma$ 与分次相容，条件 (5) 表明当 $n>1$ 时，
$\gamma_n$ 在奇次数齐次元素的乘积上为零。但请注意，仍有可能
$$
\gamma_2(z_1 z_2 + z_3 z_4) = z_1z_2z_3z_4
$$
非零，其中 $z_1,z_2,z_3,z_4$ 是奇次数齐次元素。

\begin{example}[添加奇次变量]
\label{dpa-example-adjoining-odd}
设 $R$ 是环。设 $(A,\gamma)$ 是配备上述除幂结构的严格分次交换
分次 $R$-代数，并设 $d>0$ 是奇数。在此情形可向 $A$ 添加一个
次数为 $d$ 的变量 $T$。具体地，令
$$
A\langle T \rangle = A \oplus AT
$$
其分次由 $A\langle T\rangle_m=A_m\oplus A_{m-d}T$ 给出。
我们断言 $A\langle T\rangle$ 上存在唯一的除幂结构，
与 $A$ 上给定的除幂结构相容。具体地，对
$x\in A_{even,+}$、$y\in A_{odd}$，令
$$
\gamma_n(x + yT) = \gamma_n(x) + \gamma_{n - 1}(x)yT
$$
\end{example}

\begin{example}[添加偶次变量]
\label{dpa-example-adjoining-even}
设 $R$ 是环。设 $(A,\gamma)$ 是配备上述除幂结构的严格分次交换
分次 $R$-代数，并设 $d>0$ 是偶数。在此情形可向 $A$ 添加一个
次数为 $d$ 的变量 $T$。具体地，令
$$
A\langle T \rangle = A \oplus AT \oplus AT^{(2)} \oplus AT^{(3)} \oplus \ldots
$$
其乘法由
$$
T^{(n)} T^{(m)} = \frac{(n + m)!}{n!m!} T^{(n + m)}
$$
给出，分次由
$$
A\langle T \rangle_m =
A_m \oplus A_{m - d}T \oplus A_{m - 2d}T^{(2)} \oplus \ldots
$$
给出。我们断言 $A\langle T\rangle$ 上存在唯一的除幂结构，
与 $A$ 上给定的除幂结构相容，并满足
$\gamma_n(T^{(i)})=T^{(ni)}$。为定义该除幂结构，先在
$x_i\in A_{even}$ 时令
$$
\gamma_n\left(\sum\nolimits_{i > 0} x_i T^{(i)}\right) =
\sum \prod\nolimits_{n = \sum e_i} x_i^{e_i} T^{(ie_i)}
$$
若 $x_0\in A_{even,+}$，则令
$$
\gamma_n\left(\sum\nolimits_{i \geq 0} x_i T^{(i)}\right) =
\sum\nolimits_{a + b = n}
\gamma_a(x_0)\gamma_b\left(\sum\nolimits_{i > 0} x_iT^{(i)}\right)
$$
其中 $\gamma_b$ 如上定义。
\end{example}

\begin{remark}
\label{dpa-remark-adjoining-set-of-variables}
也可以在给定次数 $d>0$ 中添加一族（可能无限多个）外代数生成元
或除幂生成元，而不只是像例 \ref{dpa-example-adjoining-odd} 和
\ref{dpa-example-adjoining-even} 那样添加一个。具体地，沿用评注
\ref{dpa-remark-divided-power-polynomial-algebra}：对如上的 $(A,\gamma)$
及集合 $J$，令 $A\langle T_j:j\in J\rangle$ 为代数
$A\langle T_j:j\in S\rangle$ 在 $J$ 的所有有限子集 $S$ 上的
有向余极限。显然，该代数具有唯一的除幂结构，并与 $A$ 上给定的
结构以及每个生成元 $T_j$ 上的结构相容。
\end{remark}

\noindent
现在把除幂结构的定义与微分联系起来。为理解下面的定义，注意若
$A$ 是 $\mathbf{Q}$-代数且 $x\in A_{even,+}$，则
$\text{d}(x^n/n!)=\text{d}(x)x^{n-1}/(n-1)!$。

\begin{definition}
\label{dpa-definition-divided-powers-dga}
设 $R$ 是环，$A=\bigoplus_{d\geq0}A_d$ 是严格分次交换的微分分次
$R$-代数。若对所有 $x\in A_{even,+}$ 有
$\text{d}(\gamma_n(x))=\text{d}(x)\gamma_{n-1}(x)$，
则称 $A$ 上的除幂结构 $\gamma$ {\it 与微分分次结构相容}。
\end{definition}

\noindent
警告：设 $(A,\text{d},\gamma)$ 如定义
\ref{dpa-definition-divided-powers-dga}。即使 $x$ 是边界，
$\gamma_n(x)$ 也未必是边界。因此，一般而言 $\gamma$ 不在同调代数
$H(A)$ 上诱导除幂结构。有些论文的作者会加入另一条相容性条件，
以保证确实诱导这样的结构；这里我们选择不这样做。

\begin{lemma}
\label{dpa-lemma-dpdga-good}
设 $(A,\text{d},\gamma)$ 和 $(B,\text{d},\gamma)$ 如定义
\ref{dpa-definition-divided-powers-dga}，并设 $f:A\to B$ 是与除幂结构
相容的微分分次代数态射。假设
\begin{enumerate}
\item 对 $k>0$，$H_k(A)=0$；
\item $f$ 是满射。
\end{enumerate}
则 $\gamma$ 在分次 $R$-代数 $H(B)$ 上诱导一个除幂结构。
\end{lemma}

\begin{proof}
设 $x$ 和 $x'$ 是次数同为 $2d$ 的齐次元素，并在 $H(B)$ 中定义
同一个上同调类。写成 $x'-x=\text{d}(w)$。选取 $x$ 的提升
$y\in A_{2d}$ 和 $w$ 的提升 $z\in A_{2d+1}$。
则 $y'=y+\text{d}(z)$ 是 $x'$ 的提升。因此
$$
\gamma_n(y') = \sum \gamma_i(y) \gamma_{n - i}(\text{d}(z))
= \gamma_n(y) +
\sum\nolimits_{i < n} \gamma_i(y) \gamma_{n - i}(\text{d}(z))
$$
由于 $A$ 在正次数无环，并且对所有 $j$ 有
$\text{d}(\gamma_j(\text{d}(z)))=0$，可将其写成
$$
\gamma_n(y') = \gamma_n(y) +
\sum\nolimits_{i < n} \gamma_i(y) \text{d}(z_i)
$$
其中 $z_i$ 是 $A$ 中的某些元素。此外，对 $0<i<n$ 有
$$
\text{d}(\gamma_i(y) z_i) =
\text{d}(\gamma_i(y))z_i + \gamma_i(y)\text{d}(z_i) =
\text{d}(y) \gamma_{i - 1}(y) z_i + \gamma_i(y)\text{d}(z_i)
$$
由于 $\text{d}(y)$ 在 $B$ 中的像为零，第一项在 $B$ 中映为零。
因此，$\gamma_n(x')$ 和 $\gamma_n(x)$ 在 $H(B)$ 中映为同一元素。
由此，对所有 $d>0$、$n>0$ 得到良定义的映射
$\gamma_n:H_{2d}(B)\to H_{2nd}(B)$。关于这些映射在 $H(B)$ 上
定义除幂结构的验证从略。
\end{proof}

\begin{lemma}
\label{dpa-lemma-base-change-div}
设 $(A,\text{d},\gamma)$ 如定义
\ref{dpa-definition-divided-powers-dga}，并设 $R\to R'$ 是环同态。
则 $\text{d}$ 和 $\gamma$ 在 $A'=A\otimes_RR'$ 上诱导同类结构，
使得 $(A',\text{d},\gamma)$ 也如定义
\ref{dpa-definition-divided-powers-dga} 所述。
\end{lemma}

\begin{proof}
注意 $A'_{even}=A_{even}\otimes_RR'$ 且
$A'_{even,+}=A_{even,+}\otimes_RR'$。因此要证明的是除幂结构
$\gamma$ 可延拓到 $A'_{even}$（术语见定义 \ref{dpa-definition-extends}）。
一旦证明 $\gamma$ 可延拓，便容易看出该延拓具有全部所需性质。

\medskip\noindent
选取一个多项式 $R$-代数 $P$（其生成元集合任意）以及一个满的
$R$-代数同态 $P\to R'$。环同态
$A_{even}\to A_{even}\otimes_RP$ 是平坦的，故由引理
\ref{dpa-lemma-gamma-extends}，除幂结构 $\gamma$ 唯一地延拓到
$A_{even}\otimes_RP$。令 $J=\Ker(P\to R')$。
为证明 $\gamma$ 可延拓到 $A\otimes_RR'$，只须证明
$I'=\Ker(A_{even,+}\otimes_RP\to A_{even,+}\otimes_RR')$
由满足 $\gamma_n(z)\in I'$（对所有 $n>0$）的元素 $z$ 生成。
这是显然的，因为 $I'$ 由形如 $x\otimes f$ 的元素生成，其中
$x\in A_{even,+}$ 且 $f\in\Ker(P\to R')$。
\end{proof}

\begin{lemma}
\label{dpa-lemma-extend-differential}
设 $(A,\text{d},\gamma)$ 如定义
\ref{dpa-definition-divided-powers-dga}，并设 $d\geq1$ 是整数。
设 $A\langle T\rangle$ 是例 \ref{dpa-example-adjoining-odd} 或
\ref{dpa-example-adjoining-even} 中构造的、以 $T$ 为变量且
$\deg(T)=d$ 的分次除幂多项式代数。设 $f\in A_{d-1}$ 且
$\text{d}(f)=0$。则 $A\langle T\rangle$ 上存在唯一的微分
$\text{d}$，使得 $\text{d}(T)=f$，并且 $\text{d}$ 与
$A\langle T\rangle$ 上的除幂结构相容。
\end{lemma}

\begin{proof}
可由直接计算证明，细节从略。
\end{proof}

\noindent
在引理 \ref{dpa-lemma-compute-cohomology-adjoin-variable} 中，
我们将在若干特殊情形计算 $A\langle T\rangle$ 的上同调。
下面给出 Tate 的构造及 Avramov 和 Halperin 对它的推广。

\begin{lemma}
\label{dpa-lemma-tate-resolution}
设 $R\to S$ 是交换环同态。存在分解
$$
R \to A \to S
$$
具有如下性质：
\begin{enumerate}
\item $(A,\text{d},\gamma)$ 如定义
\ref{dpa-definition-divided-powers-dga} 所述；
\item 若给 $S$ 配备零微分，则 $A\to S$ 是拟同构；
\item $A_0=R[x_j:j\in J]\to S$ 是任意一个从多项式环到 $S$ 的满射；
\item $A$ 是 $R$ 上的分次除幂多项式代数。
\end{enumerate}
最后一个条件是指：从 $A_0$ 出发，像例 \ref{dpa-example-adjoining-odd}
或 \ref{dpa-example-adjoining-even} 那样，在每个正次数逐次添加一族变量
$T$，从而构造出 $A$。此外，若 $R$ 是诺特环且 $R\to S$ 为有限型，
则可以选取 $A$，使其在每个次数只有有限多个生成元。
\end{lemma}

\begin{proof}
我们写出 $R$ 为诺特环且 $R\to S$ 为有限型时的构造。
没有这些假设时，证明相同，只不过每个次数须使用一族
（可能无限多个）生成元。

\medskip\noindent
构造的起点：令 $A(0)=R[x_1,\ldots,x_n]$ 为通常的多项式环，
并设 $A(0)\to S$ 是满射。取分次
$A(0)_0=A(0)$，且当 $d\not=0$ 时 $A(0)_d=0$。
因此 $\text{d}=0$，并且当 $n>0$ 时 $\gamma_n$ 也为零。

\medskip\noindent
选取给定映射 $A(0)=R[x_1,\ldots,x_n]\to S$ 的核的一组生成元
$f_1,\ldots,f_m\in R[x_1,\ldots,x_n]$。
应用例 \ref{dpa-example-adjoining-odd} 共 $m$ 次，得到
$$
A(1) = A(0)\langle T_1, \ldots, T_m\rangle
$$
作为分次除幂多项式代数，其中 $\deg(T_i)=1$。令
$\text{d}(T_i)=f_i$。由于 $A(1)$ 是 $A(0)$ 上的除幂多项式代数，
且 $\text{d}(f_i)=0$，由引理 \ref{dpa-lemma-extend-differential}，
这唯一地延拓为 $A(1)$ 上的微分。

\medskip\noindent
归纳假设：设已给出分解
$$
R \to A(0) \to A(1) \to \ldots \to A(m) \to S
$$
其中 $A(0)$、$A(1)$ 如上；对 $2\leq m'\leq m$，每个
$R\to A(m')\to S$ 都满足引理陈述中的性质 (1)、(4)，
而将 (2) 换成条件：当 $m'>i\geq0$ 时，
$H_i(A(m'))\to H_i(S)$ 是同构。基础情形是 $m=1$。

\medskip\noindent
归纳步骤：假设已有归纳假设中的 $R\to A(m)\to S$。
考虑群 $H_m(A(m))$，它是 $H_0(A(m))=S$ 上的模。
事实上，它是 $A(m)_m$ 的一个子商，而 $A(m)_m$ 是
$A(m)_0=R[x_1,\ldots,x_n]$ 上的有限型模。因此可选取有限多个元素
$$
e_1, \ldots, e_t \in \Ker(\text{d} : A(m)_m \to A(m)_{m - 1})
$$
其像生成该模。应用例 \ref{dpa-example-adjoining-odd} 或
\ref{dpa-example-adjoining-even} 共 $t$ 次，得到
$$
A(m + 1) = A(m)\langle T_1, \ldots, T_t\rangle
$$
作为分次除幂代数，其中 $\deg(T_i)=m+1$。令
$\text{d}(T_i)=e_i$。由于 $A(m+1)$ 是 $A(m)$ 上的除幂多项式代数，
且 $\text{d}(e_i)=0$，这唯一地延拓为 $A(m+1)$ 上与除幂结构相容的
微分。由于只在次数 $m+1$ 及更高次数添加了材料，故当 $i<m$ 时
$H_i(A(m+1))=H_i(A(m))$。此外，由构造显然有
$H_m(A(m+1))=0$。

\medskip\noindent
为完成证明，注意我们已经证明存在映射序列
$$
R \to A(0) \to A(1) \to \ldots \to A(m) \to A(m + 1) \to \ldots \to S
$$
令 $A=\colim A(m)$ 即完成证明。
\end{proof}

\begin{lemma}
\label{dpa-lemma-tate-resoluton-pseudo-coherent-ring-map}
设 $R\to S$ 是伪凝聚环同态（《代数详论》定义
\ref{more-algebra-definition-pseudo-coherent-perfect}）。则引理
\ref{dpa-lemma-tate-resolution} 成立，并且可使 $S$ 的分辨 $A$
在每个次数只有有限多个生成元。
\end{lemma}

\begin{proof}
证明与引理 \ref{dpa-lemma-tate-resolution} 完全相同。唯一额外之处是：
给定 $A(m)\to S$，须证明 $H_m=H_m(A(m))$ 是有限
$R[x_1,\ldots,x_m]$-模（从而下一步只须添加有限多个变量）。
考虑复形
$$
\ldots \to A(m)_{m - 1} \to A(m)_m \to A(m)_{m - 1} \to
\ldots \to A(m)_0 \to S \to 0
$$
由于 $S$ 是伪凝聚 $R[x_1,\ldots,x_n]$-模，且 $A(m)_i$ 是有限自由
$R[x_1,\ldots,x_n]$-模，故该复形是伪凝聚复形，参见
《代数详论》引理 \ref{more-algebra-lemma-complex-pseudo-coherent-modules}。
由于该复形在同调次数 $>m$ 正合，由《代数详论》引理
\ref{more-algebra-lemma-finite-cohomology}，$H_m$ 是有限 $R$-模。
\end{proof}

\begin{lemma}
\label{dpa-lemma-uniqueness-tate-resolution}
设 $R$ 是交换环。假设 $(A,\text{d},\gamma)$ 和
$(B,\text{d},\gamma)$ 如定义 \ref{dpa-definition-divided-powers-dga}。
设 $\overline{\varphi}:H_0(A)\to H_0(B)$ 是 $R$-代数同态。假设
\begin{enumerate}
\item $A$ 是 $R$ 上的分次除幂多项式代数；
\item 对 $k>0$，$H_k(B)=0$。
\end{enumerate}
则存在与除幂结构相容的微分分次 $R$-代数态射
$\varphi:A\to B$，它提升 $\overline{\varphi}$。
\end{lemma}

\begin{proof}
该假设意味着：从 $R$ 出发，依次添加次数为零的某族多项式生成元、
正奇次数的外代数生成元和正偶次数的除幂生成元，便得到 $A$。
因此 $A$ 具有滤过 $R\subset A(0)\subset A(1)\subset\ldots$，
其中 $A(m+1)$ 由 $A(m)$ 添加次数为 $m+1$ 的适当类型生成元而得
（我们将它们统称为“除幂生成元”）。特别地，$A(0)\to H_0(A)$
是从 $R$ 上通常的多项式代数到 $H_0(A)$ 的满射。因此可将
$\overline{\varphi}$ 提升为 $R$-代数同态
$\varphi(0):A(0)\to B_0$。

\medskip\noindent
将 $A(1)$ 写成 $A(0)\langle T_j:j\in J\rangle$，其中 $J$ 是一族
次数为 $1$ 的除幂变量 $T_j$ 的指标集。令
$f_j=\varphi(0)(\text{d}(T_j))\in B_0$。由于 $\text{d}T_j$
在 $H_0(A)$ 中映为零，$f_j$ 在 $H_0(B)$ 中也映为零。
因此可找到 $b_j\in B_1$，使得 $\text{d}(b_j)=f_j$。
由引理 \ref{dpa-lemma-divided-power-polynomial-algebra} 中除幂多项式代数的
泛性质，可找到 $\varphi(0)$ 的提升 $\varphi(1):A(1)\to B$，
将 $T_j$ 映为 $f_j$。

\medskip\noindent
对某个 $m\geq1$ 构造出 $\varphi(m)$ 后，可完全以相同方式构造
$\varphi(m+1):A(m+1)\to B$。细节从略。
\end{proof}

\begin{lemma}
\label{dpa-lemma-divided-powers-on-tor}
设 $R$ 是交换环，$S$、$T$ 是交换 $R$-代数。则
$$
\operatorname{Tor}_*^R(S, T).
$$
上具有典范的、配备除幂结构的严格分次交换 $R$-代数结构。
\end{lemma}

\begin{proof}
选取如上的分解 $R\to A\to S$。由于 $A\to S$ 是拟同构，且
$A_d$ 是自由 $R$-模，微分分次代数 $B=A\otimes_RT$ 计算引理中
所示的 Tor 群。选取满射 $R[y_j:j\in J]\to T$。
于是 $B$ 是微分分次代数 $A[y_j:j\in J]$ 的商；后者的同调集中在
次数 $0$（等于 $S[y_j:j\in J]$）。由引理
\ref{dpa-lemma-base-change-div}，微分分次代数 $B$ 和
$A[y_j:j\in J]$ 都具有与微分相容的除幂结构。因此，由引理
\ref{dpa-lemma-dpdga-good}，在 $H(B)$ 上得到所需的除幂结构。

\medskip\noindent
以此方式构造的除幂代数结构与 $A$ 的选择无关。事实上，若 $A'$ 是
另一种选择，则引理 \ref{dpa-lemma-uniqueness-tate-resolution} 蕴含存在
保持全部结构以及到 $S$ 的增广的态射 $A\to A'$。
于是诱导态射
$B=A\otimes_RT\to A'\otimes_RT'=B'$ 也保持全部结构并且是拟同构。
因此，诱导的 Tor 代数同构与乘积和除幂相容。
\end{proof}





\section{在完全交中的应用}
\label{dpa-section-application-ci}

\noindent
设 $R$ 是环，$(A,\text{d},\gamma)$ 如定义
\ref{dpa-definition-divided-powers-dga}。次数为 $2$ 的{\it 导子}是具有
下列性质的 $R$-线性映射 $\theta:A\to A$：
\begin{enumerate}
\item $\theta(A_d) \subset A_{d - 2}$,
\item $\theta(xy) = \theta(x)y + x\theta(y)$,
\item $\theta$ 与 $\text{d}$ 交换；
\item $\theta(\gamma_n(x)) = \theta(x) \gamma_{n - 1}(x)$
对所有 $d$ 及所有 $x\in A_{2d}$ 成立。
\end{enumerate}
下述引理构造一个导子。

\begin{lemma}
\label{dpa-lemma-get-derivation}
设 $R$ 是环，$(A,\text{d},\gamma)$ 如定义
\ref{dpa-definition-divided-powers-dga}。设 $R'\to R$ 是环的满射，
其核平方为零且由一个元素 $f$ 生成。若 $A$ 是 $R$ 上的分次
除幂多项式代数，并且每个次数只有有限多个变量，则得到导子
$\theta:A/IA\to A/IA$，其中 $I$ 是 $f$ 在 $R$ 中的零化子。
\end{lemma}

\begin{proof}
由于 $A$ 是除幂多项式代数，可找到 $R'$ 上的除幂多项式代数
$A'$，使得 $A=A'\otimes_{R'}R$。此外，可以把 $\text{d}$ 提升为
$A'$ 上的 $R$-线性算子 $\text{d}$，使得
\begin{enumerate}
\item  $\text{d}(xy) = \text{d}(x)y + (-1)^{\deg(x)}x \text{d}(y)$
对 $A'$ 中的齐次元素 $x,y$ 成立；
\item  $\text{d}(\gamma_n(x)) = \text{d}(x) \gamma_{n - 1}(x)$ for
$x \in A'_{even, +}$。
\end{enumerate}
细节从略（提示：每次处理一个变量）。不过，$\text{d}^2$ 在 $A'$
上未必为零。显然，$\text{d}^2$ 将 $A'$ 映入
$fA'\cong A/IA$。因此 $\text{d}^2$ 零化 $fA'$，并分解为映射
$A\to A/IA$。由于 $\text{d}^2$ 是 $R$-线性的，得到所需映射
$\theta:A/IA\to A/IA$。导子各项性质的验证是立即可得的。
\end{proof}

\begin{lemma}
\label{dpa-lemma-compute-theta}
假设和记号如引理 \ref{dpa-lemma-get-derivation}。假设 $S=H_0(A)$ 同构于
$R[x_1, \ldots, x_n]/(f_1, \ldots, f_m)$
，其中 $n,m$ 为某些整数且 $f_j\in R[x_1,\ldots,x_n]$。
再假设给定关系
$$
\sum r_j f_j = 0
$$
其中 $r_j\in R[x_1,\ldots,x_n]$。选取 $r_j,f_j$ 的提升
$r'_j,f'_j\in R'[x_1,\ldots,x_n]$。将
$\sum r'_jf'_j$ 写成 $gf$，其中 $g\in R/I[x_1,\ldots,x_n]$。
若 $H_1(A)=0$ 且每个 $r_j$ 的所有系数都属于 $I$，则存在元素
$\xi\in H_2(A/IA)$，使得在 $S/IS$ 中 $\theta(\xi)=g$。
\end{lemma}

\begin{proof}
设 $A(0)\subset A(1)\subset A(2)\subset\ldots$ 是 $A$ 的滤过，
使得 $A(m)$ 由 $A(m-1)$ 添加次数为 $m$ 的除幂变量而得。
则 $A(0)$ 是 $R$ 上的多项式代数，并配备一个 $R$-代数满射
$A(0)\to S$。因此可以选取提升到 $S$ 的增广的映射
$$
\varphi : R[x_1, \ldots, x_n] \to A(0)
$$
接着，对某些次数为 $1$ 的除幂变量 $T_i$，有
$A(1)=A(0)\langle T_1,\ldots,T_t\rangle$。由于 $H_0(A)=S$，
可选取 $\xi_j\in\sum A(0)T_i$，使得
$\text{d}(\xi_j)=\varphi(f_j)$。于是
$$
\text{d}\left(\sum \varphi(r_j) \xi_j\right) =
\sum  \varphi(r_j) \varphi(f_j) =  \sum \varphi(r_jf_j) = 0
$$
由于 $H_1(A)=0$，可选取 $\xi\in A_2$，使得
$\text{d}(\xi)=\sum\varphi(r_j)\xi_j$。若 $r_j$ 的系数属于 $I$，
则 $\varphi(r_j)$ 也如此。此时 $\text{d}(\xi)$ 在 $A_1/IA_1$
中消失，因而 $\xi$ 在 $H_2(A/IA)$ 中定义一个类。

\medskip\noindent
引理 \ref{dpa-lemma-get-derivation} 的证明中对 $\theta$ 的构造，是逐次将
$A(i)$ 提升到 $A'(i)$ 并提升微分 $\text{d}$。将 $\varphi$ 提升为
$\varphi':R'[x_1,\ldots,x_n]\to A'(0)$。接着，
$A'(1)=A'(0)\langle T_1,\ldots,T_t\rangle$。此外，可将 $\xi_j$
提升为 $\xi'_j\in\sum A'(0)T_i$。于是，对某个 $a_j\in A'(0)$，
$\text{d}(\xi'_j)=\varphi'(f'_j)+fa_j$。
考虑 $\xi$ 的一个提升 $\xi'\in A'_2$。则
$$
\text{d}(\xi') = \sum \varphi'(r'_j)\xi'_j + \sum fb_iT_i
$$
其中 $b_i\in A(0)$。再次应用 $\text{d}$，得到
$$
\theta(\xi) = \sum \varphi'(r'_j)\varphi'(f'_j) +
\sum f \varphi'(r'_j) a_j + \sum fb_i \text{d}(T_i)
$$
第一项给出所需结果。第二项为零，因为 $r_j$ 的系数属于 $I$，
从而被 $f$ 零化。第三项在 $H_0$ 中映为零，因为
$\text{d}(T_i)$ 映为零。
\end{proof}

\noindent
下述引理的证明方法似乎源自 Gulliksen。

\begin{lemma}
\label{dpa-lemma-not-finite-pd}
设 $R'\to R$ 是诺特环的满射，其核平方为零且由一个元素 $f$ 生成。令
$S = R[x_1, \ldots, x_n]/(f_1, \ldots, f_m)$.
设 $\sum r_jf_j=0$ 是 $R[x_1,\ldots,x_n]$ 中的一个关系。假设
\begin{enumerate}
\item 每个 $r_j$ 的系数都属于 $f$ 在 $R$ 中的零化子 $I$；
\item 对某些提升 $r'_j,f'_j\in R'[x_1,\ldots,x_n]$，有
$\sum r'_jf'_j=gf$，其中 $g$ 在 $S/IS$ 中不幂零。
\end{enumerate}
则 $S$ 在 $R$ 上不具有有限 Tor 维数（即 $S$ 不是完美 $R$-代数）。
\end{lemma}

\begin{proof}
选取引理 \ref{dpa-lemma-tate-resolution} 中的 Tate 分辨
$R\to A\to S$。设 $\xi\in H_2(A/IA)$ 和
$\theta:A/IA\to A/IA$ 是引理 \ref{dpa-lemma-get-derivation} 和
\ref{dpa-lemma-compute-theta} 中得到的元素和导子。注意
$$
\theta^n(\gamma_n(\xi)) = g^n
$$
在 $H_0(A/IA)=S/IS$ 中成立。因此，若 $g$ 在 $S/IS$ 中不幂零，
则对所有 $n>0$，$\xi^n$ 在 $H_{2n}(A/IA)$ 中非零。由于
$H_{2n}(A/IA)=\text{Tor}^R_{2n}(S,R/I)$，结论成立。
\end{proof}

\noindent
下述结果可见于 \cite{Rodicio}。

\begin{lemma}
\label{dpa-lemma-injective}
设 $(A,\mathfrak m)$ 是诺特局部环，$I\subset J\subset A$ 是真理想。
若 $A/J$ 在 $A/I$ 上具有有限 Tor 维数，则
$I/\mathfrak mI\to J/\mathfrak mJ$ 是单射。
\end{lemma}

\begin{proof}
设 $f\in I$ 在 $I/\mathfrak mI$ 中的像非零，而在
$J/\mathfrak mJ$ 中映为零。可以选取理想 $I'$，满足
$\mathfrak mI\subset I'\subset I$，使得 $I/I'$ 由 $f$ 的像生成。
令 $R=A/I$、$R'=A/I'$。设 $J=(a_1,\ldots,a_m)$，其中
$a_j\in A$。则对某些 $b_j\in\mathfrak m$ 有
$f=\sum b_ja_j$。分别令 $r_j,f_j\in R$ 和
$r'_j,f'_j\in R'$ 为 $b_j,a_j$ 的像。于是处于引理
\ref{dpa-lemma-not-finite-pd} 的情形（其中该引理的理想 $I$ 等于
$\mathfrak m_R$），引理得证。
\end{proof}

\begin{lemma}
\label{dpa-lemma-regular-sequence}
设 $(A,\mathfrak m)$ 是诺特局部环，$I\subset J\subset A$ 是真理想。
假设
\begin{enumerate}
\item $A/J$ 在 $A/I$ 上具有有限 Tor 维数；
\item $J$ 由正则序列生成。
\end{enumerate}
则 $I$ 由正则序列生成，且 $J/I$ 由正则序列生成。
\end{lemma}

\begin{proof}
由引理 \ref{dpa-lemma-injective}，映射
$I/\mathfrak mI\to J/\mathfrak mJ$ 是单射。因此可找到 $s\leq r$
以及 $J$ 的一个极小生成元系 $f_1,\ldots,f_r$，使得
$f_1,\ldots,f_s$ 属于 $I$ 并构成 $I$ 的极小生成元系。
由《代数详论》引理
\ref{more-algebra-lemma-independence-of-generators} 和
\ref{more-algebra-lemma-noetherian-finite-all-equivalent}.
，$J$ 的任意极小生成元系都是正则序列，故引理成立。
\end{proof}

\begin{lemma}
\label{dpa-lemma-perfect-map-ci}
设 $R\to S$ 是诺特局部环之间的局部环同态。
设 $I\subset R$、$J\subset S$ 是理想且 $IS\subset J$。
若 $R\to S$ 平坦且 $S/\mathfrak m_RS$ 正则，则下列条件等价：
\begin{enumerate}
\item $J$ 由正则序列生成，且 $S/J$ 作为 $R/I$-模具有有限 Tor 维数；
\item $J$ 由正则序列生成，且只有有限多个 $p$ 使
$\text{Tor}^{R/I}_p(S/J,R/\mathfrak m_R)$ 非零；
\item $I$ 由正则序列生成，且 $J/IS$ 在 $S/IS$ 中由正则序列生成。
\end{enumerate}
\end{lemma}

\begin{proof}
若 (3) 成立，则 $J$ 由正则序列生成，例如参见《代数详论》引理
\ref{more-algebra-lemma-join-koszul-regular-sequences} 和
\ref{more-algebra-lemma-noetherian-finite-all-equivalent}.
此外，若 (3) 成立，则 $S/J=(S/I)/(J/I)$ 在 $S/IS$ 上具有有限
投射维数，因为 Koszul 复形是 $S/J$ 在 $S/IS$ 上的有限自由分辨。
由于 $R/I\to S/IS$ 平坦，由《代数详论》引理
\ref{more-algebra-lemma-flat-push-tor-amplitude}，$S/J$ 在 $R/I$
上具有有限 Tor 维数。因此 (3) 蕴含 (1)。

\medskip\noindent
(1) $\Rightarrow$ (2) 是平凡的。假设 (2)。由《代数详论》引理
\ref{more-algebra-lemma-perfect-over-regular-local-ring}，$S/J$ 在
$S/IS$ 上具有有限 Tor 维数。因此可应用引理
\ref{dpa-lemma-regular-sequence}，得到 $IS$ 和 $J/IS$ 都由正则序列生成。
设 $f_1,\ldots,f_r\in I$ 是 $I$ 的极小生成元系。由于 $R\to S$
平坦，$f_1,\ldots,f_r$ 在 $S$ 中构成 $IS$ 的极小生成元系。
因此，由《代数详论》引理
\ref{more-algebra-lemma-independence-of-generators} 和
\ref{more-algebra-lemma-noetherian-finite-all-equivalent}.
，元素序列 $f_1,\ldots,f_r\in R$ 的像在 $S$ 中构成正则序列。
故由《代数》引理 \ref{algebra-lemma-flat-increases-depth}，
$f_1,\ldots,f_r$ 在 $R$ 中是正则序列。
\end{proof}





\section{局部完全交环}
\label{dpa-section-lci}

\noindent
设 $(A,\mathfrak m)$ 是诺特完备局部环。由 Cohen 结构定理
（参见《代数》定理 \ref{algebra-theorem-cohen-structure-theorem}），
可将 $A$ 写成某个正则诺特完备局部环 $R$ 的商。
若存在某个满射 $R\to A$，其中 $R$ 是正则局部环且其核由正则序列
生成，则称 $A$ 是{\it 完全交}。下述引理表明该概念与满射的选择无关。

\begin{lemma}
\label{dpa-lemma-ci-well-defined}
设 $(A,\mathfrak m)$ 是诺特完备局部环。下列条件等价：
\begin{enumerate}
\item 对每个局部环满射 $R\to A$，只要 $R$ 是正则局部环，
$R\to A$ 的核便由正则序列生成；
\item 对某个局部环满射 $R\to A$，$R$ 是正则局部环，
且 $R\to A$ 的核由正则序列生成。
\end{enumerate}
\end{lemma}

\begin{proof}
设 $k$ 是 $A$ 的剩余域。若 $k$ 的特征为 $p>0$，则用 $\Lambda$
表示一个剩余域为 $k$ 的 Cohen 环（《代数》定义
\ref{algebra-definition-cohen-ring} 和引理
\ref{algebra-lemma-cohen-rings-exist}）。若 $k$ 的特征为 $0$，
令 $\Lambda=k$。回顾：对任意 $n$，
$\Lambda[[x_1,\ldots,x_n]]$ 在 $\mathfrak m$-进拓扑下分别在
$\mathbf{Z}$ 或 $\mathbf{Q}$ 上形式光滑，参见《代数详论》引理
\ref{more-algebra-lemma-power-series-ring-over-Cohen-fs}。
固定 Cohen 结构定理（《代数》定理
\ref{algebra-theorem-cohen-structure-theorem}）中的一个满射
$\Lambda[[x_1,\ldots,x_n]]\to A$。

\medskip\noindent
设 $R\to A$ 是从正则局部环 $R$ 出发的满射，并设
$f_1,\ldots,f_r$ 是 $\Ker(R\to A)$ 的极小生成元序列。
以下将反复使用而不再说明：诺特局部环中的理想由正则序列生成，
当且仅当其任意极小生成元集都是正则序列。由《代数》引理
\ref{algebra-lemma-flat-increases-depth} 和
\ref{algebra-lemma-completion-flat}，$f_1,\ldots,f_r$ 在 $R$ 中为
正则序列，当且仅当它们在完备化 $R^\wedge$ 中为正则序列。此外，
$$
R^\wedge/(f_1, \ldots, f_r)R^\wedge =
(R/(f_1, \ldots, f_n))^\wedge = A^\wedge = A
$$
因为 $A$ 是 $\mathfrak m_A$-进完备的（第一个等式由《代数》引理
\ref{algebra-lemma-completion-tensor} 得到）。最后，由于 $R$ 正则，
$R^\wedge$ 也正则（《代数详论》引理
\ref{more-algebra-lemma-completion-regular}）。故可假设 $R$ 完备。

\medskip\noindent
若 $R$ 完备，则可选取映射
$\Lambda[[x_1,\ldots,x_n]]\to R$，提升给定映射
$\Lambda[[x_1,\ldots,x_n]]\to A$，参见《代数详论》引理
\ref{more-algebra-lemma-lift-continuous}。再添加一些变量
$y_1,\ldots,y_m$，将其映到 $R\to A$ 的核的生成元，便可假设
$\Lambda[[x_1,\ldots,x_n,y_1,\ldots,y_m]]\to R$ 是满射
（略去若干细节）。于是可考虑交换图
$$
\xymatrix{
\Lambda[[x_1, \ldots, x_n, y_1, \ldots, y_m]] \ar[r] \ar[d] & R \ar[d] \\
\Lambda[[x_1, \ldots, x_n]] \ar[r] & A
}
$$
由《代数》引理 \ref{algebra-lemma-ci-well-defined}，关于 $R\to A$
的条件等价于关于固定映射
$\Lambda[[x_1,\ldots,x_n]]\to A$ 的条件。引理得证。
\end{proof}

\noindent
下面两个引理是对上述定义的合理性检验。

\begin{lemma}
\label{dpa-lemma-quotient-regular-ring-by-regular-sequence}
设 $R$ 是正则环，$\mathfrak p\subset R$ 是素理想，且
$f_1,\ldots,f_r\in\mathfrak p$ 是正则序列。则
$$
A = (R/(f_1, \ldots, f_r))_\mathfrak p =
R_\mathfrak p/(f_1, \ldots, f_r)R_\mathfrak p
$$
的完备化是上述意义下的完全交。
\end{lemma}

\begin{proof}
由诺特环 $R_\mathfrak p$ 上有限模的完备化是正合的
（《代数》引理 \ref{algebra-lemma-completion-tensor}），$A$ 的完备化为
$A^\wedge = R_\mathfrak p^\wedge/(f_1, \ldots, f_r)R_\mathfrak p^\wedge$
。由《代数》引理 \ref{algebra-lemma-completion-flat} 和
\ref{algebra-lemma-flat-increases-depth}，序列 $f_1,\ldots,f_r$ 在
$R_\mathfrak p$ 中的像是正则序列。此外，由《代数详论》引理
\ref{more-algebra-lemma-completion-regular}，$R_\mathfrak p^\wedge$
是正则局部环。因此结论由完备局部环的完全交定义成立。
\end{proof}

\noindent
下述引理是《代数》引理 \ref{algebra-lemma-lci} 的对应版本。

\begin{lemma}
\label{dpa-lemma-quotient-regular-ring}
设 $R$ 是正则环，$\mathfrak p\subset R$ 是素理想，
$I\subset\mathfrak p$ 是理想。令
$A=(R/I)_\mathfrak p=R_\mathfrak p/I_\mathfrak p$。下列条件等价：
\begin{enumerate}
\item $A$ 的完备化是上述意义下的完全交；
\item $I_\mathfrak p\subset R_\mathfrak p$ 由正则序列生成；
\item 模 $(I/I^2)_\mathfrak p$ 可由
$\dim(R_\mathfrak p)-\dim(A)$ 个元素生成；
\item 此处再添加内容。
\end{enumerate}
\end{lemma}

\begin{proof}
可以并确实将 $R$ 替换为它在 $\mathfrak p$ 处的局部化。
于是 $\mathfrak p=\mathfrak m$ 是 $R$ 的极大理想，且 $A=R/I$。
设 $f_1,\ldots,f_r\in I$ 是极小生成元序列。由诺特环
$R_\mathfrak p$ 上有限模的完备化是正合的（《代数》引理
\ref{algebra-lemma-completion-tensor}），$A$ 的完备化为
$A^\wedge = R^\wedge/(f_1, \ldots, f_r)R^\wedge$
。

\medskip\noindent
若 (1) 成立，则由假设，序列 $f_1,\ldots,f_r$ 在 $R^\wedge$ 中的像
是正则序列。故由《代数》引理 \ref{algebra-lemma-completion-flat} 和
\ref{algebra-lemma-flat-increases-depth}，它在 $R$ 中也是正则序列。
因此 (1) 蕴含 (2)。

\medskip\noindent
假设 (3) 成立。令 $c=\dim(R)-\dim(A)$，并取
$f_1,\ldots,f_c\in I$，使其像生成 $I/I^2$。由 Nakayama 引理
（《代数》引理 \ref{algebra-lemma-NAK}），
$I=(f_1,\ldots,f_c)$。由于 $R$ 正则，因而是 Cohen--Macaulay 的
（《代数》命题 \ref{algebra-proposition-CM-module}），再由《代数》命题
\ref{algebra-proposition-CM-module}，$f_1,\ldots,f_c$ 是正则序列。
因此 (3) 蕴含 (2)。最后，由引理
\ref{dpa-lemma-quotient-regular-ring-by-regular-sequence}，(2) 蕴含 (1)。
\end{proof}

\noindent
下述结果由 Avramov 得到，参见 \cite{Avramov}。

\begin{proposition}
\label{dpa-proposition-avramov}
设 $A\to B$ 是诺特局部环之间的平坦局部同态。下列条件等价：
\begin{enumerate}
\item $B^\wedge$ 是完全交；
\item $A^\wedge$ 和 $(B/\mathfrak m_AB)^\wedge$ 都是完全交。
\end{enumerate}
\end{proposition}

\begin{proof}
考虑交换图
$$
\xymatrix{
B \ar[r] & B^\wedge \\
A \ar[u] \ar[r] & A^\wedge \ar[u]
}
$$
由于水平映射忠实平坦（《代数》引理
\ref{algebra-lemma-completion-faithfully-flat}），右侧竖直箭头平坦
（例如由《代数》引理
\ref{algebra-lemma-criterion-flatness-fibre-Noetherian}）。此外，
$(B/\mathfrak m_A B)^\wedge = B^\wedge/\mathfrak m_{A^\wedge} B^\wedge$
由《代数》引理 \ref{algebra-lemma-completion-tensor} 成立。
故可假设 $A$、$B$ 是完备诺特局部环。

\medskip\noindent
假设 $A$、$B$ 是完备诺特局部环。选取《代数详论》引理
\ref{more-algebra-lemma-embed-map-Noetherian-complete-local-rings}
中的交换图
$$
\xymatrix{
S \ar[r] & B \\
R \ar[u] \ar[r] & A \ar[u]
}
$$
令 $I=\Ker(R\to A)$、$J=\Ker(S\to B)$。注意，由于
$R/I=A\to B=S/J$ 平坦，映射
$J/IS \otimes_R R/\mathfrak m_R \to J/J \cap \mathfrak m_R S$
是同构。因此，$J/IS$ 的一个极小生成元系映到
$\Ker(S/\mathfrak m_RS\to B/\mathfrak m_AB)$ 的一个极小生成元系。
最后，$S/\mathfrak m_RS$ 是正则局部环。

\medskip\noindent
假设 (1) 成立，即 $J$ 由正则序列生成。由于
$A=R/I\to B=S/J$ 平坦，引理 \ref{dpa-lemma-perfect-map-ci} 适用，
从而 $I$ 和 $J/IS$ 都由正则序列生成。由《代数》引理
\ref{algebra-lemma-dimension-base-fibre-equals-total}，有
$\dim(B)=\dim(A)+\dim(B/\mathfrak m_AB)$ 且
$\dim(S/IS)=\dim(A)+\dim(S/\mathfrak m_RS)$。
因此，$J/IS$ 由
$$
\dim(S/IS) - \dim(S/J) = \dim(S/\mathfrak m_R S) - \dim(B/\mathfrak m_A B)
$$
个元素生成（《代数》引理 \ref{algebra-lemma-one-equation}）。
因此，$\Ker(S/\mathfrak m_RS\to B/\mathfrak m_AB)$ 也由同样数目的
元素生成（见上）。故该核由正则序列生成，例如参见引理
\ref{dpa-lemma-quotient-regular-ring}。由此得到 (2)。

\medskip\noindent
若 (2) 成立，则 $I$ 和 $J/J\cap\mathfrak m_RS$ 都由正则序列生成。
提升这些生成元（见上），利用 $R/I\to S/IS$ 的平坦性以及
Grothendieck 引理（《代数》引理
\ref{algebra-lemma-grothendieck-regular-sequence}），可知 $J/IS$
在 $S/IS$ 中由正则序列生成。于是引理
\ref{dpa-lemma-perfect-map-ci} 表明 $J$ 由正则序列生成，故 (1) 成立。
\end{proof}

\begin{definition}
\label{dpa-definition-lci}
设 $A$ 是诺特环。
\begin{enumerate}
\item 若 $A$ 是局部环，并且其完备化是上述意义下的完全交，
则称 $A$ 是{\it 完全交}。
\item 一般而言，若 $A$ 的所有局部环都是完全交，则称 $A$
是{\it 局部完全交}。
\end{enumerate}
\end{definition}

\noindent
下面将检验这与《代数》定义 \ref{algebra-definition-lci-field} 和
\ref{algebra-definition-lci-local-ring} 中引入的术语不冲突。
但首先说明该定义“合理”：若 $A$ 是诺特局部完全交，则 $A$
是局部完全交，即它的所有局部环都是完全交。

\begin{lemma}
\label{dpa-lemma-ci-good}
设 $(A,\mathfrak m)$ 是诺特局部环，$\mathfrak p\subset A$ 是素理想。
若 $A$ 是完全交，则 $A_\mathfrak p$ 也是完全交。
\end{lemma}

\begin{proof}
选取 $A^\wedge$ 中位于 $\mathfrak p$ 上方的素理想 $\mathfrak q$
（这是可能的，因为由《代数》引理
\ref{algebra-lemma-completion-faithfully-flat}，$A\to A^\wedge$
忠实平坦）。于是 $A_\mathfrak p\to(A^\wedge)_\mathfrak q$
是平坦局部环同态。因此，由命题 \ref{dpa-proposition-avramov}，
$A_\mathfrak p$ 是完全交，当且仅当 $(A^\wedge)_\mathfrak q$
是完全交。故只须在 $A$ 完备时证明该引理（这是证明的关键步骤）。

\medskip\noindent
假设 $A$ 完备。由定义，可将 $A$ 写成
$R/(f_1,\ldots,f_r)$，其中 $f_1,\ldots,f_r$ 是正则局部环 $R$
中的正则序列。设 $\mathfrak q\subset R$ 是对应于 $\mathfrak p$
的素理想。注意 $f_1,\ldots,f_r\in\mathfrak q$，且
$A_\mathfrak p=R_\mathfrak q/(f_1,\ldots,f_r)R_\mathfrak q$。
故由引理 \ref{dpa-lemma-quotient-regular-ring-by-regular-sequence}，
$A_\mathfrak p$ 是完全交。
\end{proof}

\begin{lemma}
\label{dpa-lemma-check-lci-at-maximal-ideals}
设 $A$ 是诺特环。则 $A$ 是局部完全交，当且仅当对 $A$ 的每个极大理想
$\mathfrak m$，$A_\mathfrak m$ 都是完全交。
\end{lemma}

\begin{proof}
立即由引理 \ref{dpa-lemma-ci-good} 和定义得到。
\end{proof}

\begin{lemma}
\label{dpa-lemma-check-lci-agrees}
设 $S$ 是域 $k$ 上的有限型代数。
\begin{enumerate}
\item 对素理想 $\mathfrak q\subset S$，局部环 $S_\mathfrak q$
是《代数》定义 \ref{algebra-definition-lci-local-ring} 意义下的完全交，
当且仅当它是定义 \ref{dpa-definition-lci} 意义下的完全交；
\item $S$ 是《代数》定义 \ref{algebra-definition-lci-field} 意义下的
局部完全交，当且仅当它是定义 \ref{dpa-definition-lci} 意义下的
局部完全交。
\end{enumerate}
\end{lemma}

\begin{proof}
证明 (1)。设 $k[x_1,\ldots,x_n]\to S$ 是满射，
$\mathfrak p\subset k[x_1,\ldots,x_n]$ 是对应于 $\mathfrak q$
的素理想，并设 $I\subset k[x_1,\ldots,x_n]$ 是该满射的核。
注意 $k[x_1,\ldots,x_n]_\mathfrak p\to S_\mathfrak q$ 是满射，
其核为 $I_\mathfrak p$。由《代数》命题
\ref{algebra-proposition-finite-gl-dim-polynomial-ring}，
$k[x_1,\ldots,x_n]$ 是正则环。因此，将引理
\ref{dpa-lemma-quotient-regular-ring} 与《代数》引理
\ref{algebra-lemma-lci-local} 结合，便得到 (1) 中两个概念的等价性。

\medskip\noindent
证明 (1) 后，(2) 中的等价性由定义和《代数》引理
\ref{algebra-lemma-lci-global} 得到。
\end{proof}

\begin{lemma}
\label{dpa-lemma-avramov}
设 $A\to B$ 是诺特局部环之间的平坦局部同态。下列条件等价：
\begin{enumerate}
\item $B$ 是完全交；
\item $A$ 和 $B/\mathfrak m_AB$ 都是完全交。
\end{enumerate}
\end{lemma}

\begin{proof}
既然定义已经成立，这只是（非平凡的）命题
\ref{dpa-proposition-avramov} 的平凡改述。
\end{proof}









\section{局部完全交同态}
\label{dpa-section-lci-homomorphisms}

\noindent
设 $A\to B$ 是诺特完备局部环之间的局部同态。Cohen 结构定理的一个
推论是，可以找到诺特完备局部环的交换图
$$
\xymatrix{
S \ar[r] & B \\
& A \ar[lu] \ar[u]
}
$$
，其中 $S\to B$ 为满射，$A\to S$ 平坦，且
$S/\mathfrak m_AS$ 是正则局部环。这由《代数详论》引理
\ref{more-algebra-lemma-embed-map-Noetherian-complete-local-rings} 得到。
暂且称 $A\to S\to B$ 是 $A\to B$ 的{\it 良好分解}，若 $S$
是诺特局部环，$A\to S\to B$ 是局部环同态，$S\to B$ 为满射，
$A\to S$ 平坦，且 $S/\mathfrak m_AS$ 正则。若存在某个良好分解
$A\to S\to B$，使得 $S\to B$ 的核由正则序列生成，则称
$A\to B$ 是{\it 完全交同态}。下述引理表明该概念与图的选择无关。

\begin{lemma}
\label{dpa-lemma-ci-map-well-defined}
设 $A\to B$ 是诺特完备局部环之间的局部同态。下列条件等价：
\begin{enumerate}
\item 对某个良好分解 $A\to S\to B$，$S\to B$ 的核由正则序列生成；
\item 对每个良好分解 $A\to S\to B$，$S\to B$ 的核都由正则序列生成。
\end{enumerate}
\end{lemma}

\begin{proof}
设 $A\to S\to B$ 是良好分解。由于 $B$ 完备，得到分解
$A\to S^\wedge\to B$，其中 $S^\wedge$ 是 $S$ 的完备化。
注意这也是良好分解：环同态 $S\to S^\wedge$ 平坦
（《代数》引理 \ref{algebra-lemma-completion-flat}），故
$A\to S^\wedge$ 平坦。由于 $S/\mathfrak m_AS$ 正则，环
$S^\wedge/\mathfrak m_AS^\wedge=(S/\mathfrak m_AS)^\wedge$ 也正则
（《代数详论》引理 \ref{more-algebra-lemma-completion-regular}）。
设 $f_1,\ldots,f_r$ 是 $\Ker(S\to B)$ 的极小生成元序列。
以下将反复使用而不再说明：诺特局部环中的理想由正则序列生成，
当且仅当其任意极小生成元集都是正则序列。由《代数》引理
\ref{algebra-lemma-flat-increases-depth}，$f_1,\ldots,f_r$ 在 $S$
中为正则序列，当且仅当它们在完备化 $S^\wedge$ 中为正则序列。
此外，
$$
S^\wedge/(f_1, \ldots, f_r)R^\wedge =
(S/(f_1, \ldots, f_n))^\wedge = B^\wedge = B
$$
因为 $B$ 是 $\mathfrak m_B$-进完备的（第一个等式由《代数》引理
\ref{algebra-lemma-completion-tensor} 得到）。因此，$S\to B$ 的核
由正则序列生成，当且仅当 $S^\wedge\to B$ 的核由正则序列生成。
故只须考虑 $S$ 完备的良好分解。

\medskip\noindent
假设有两个分解 $A\to S\to B$ 和 $A\to S'\to B$，其中 $S$、$S'$
完备。由《代数详论》引理
\ref{more-algebra-lemma-dominate-two-surjections}，环
$S\times_BS'$ 是诺特完备局部环。因此，利用《代数详论》引理
\ref{more-algebra-lemma-embed-map-Noetherian-complete-local-rings}，
可选取良好分解 $A\to S''\to S\times_BS'$，其中 $S''$ 完备。
故只须证明：若 $A\to S'\to S\to B$ 是可比较的良好分解，则
$\Ker(S\to B)$ 由正则序列生成，当且仅当 $\Ker(S'\to B)$
由正则序列生成。

\medskip\noindent
设 $A\to S'\to S\to B$ 是可比较的良好分解。首先，由于
$S'/\mathfrak m_RS'\to S/\mathfrak m_RS$ 是正则局部环的满射，
其核由正则序列
$\overline{x}_1, \ldots, \overline{x}_c \in
\mathfrak m_{S'}/\mathfrak m_R S'$
生成；该序列可扩充为正则局部环 $S'/\mathfrak m_RS'$ 的一个
正则参数系，参见《代数》引理
\ref{algebra-lemma-regular-quotient-regular}。令 $I=\Ker(S'\to S)$。
由 $S$ 在 $R$ 上平坦，有
$$
I/\mathfrak m_R I =
\Ker(S'/\mathfrak m_R S' \to S/\mathfrak m_R S) =
(\overline{x}_1, \ldots, \overline{x}_c).
$$
选取提升 $x_1,\ldots,x_c\in I$。由于 $S'$ 在 $R$ 上平坦，
这些提升构成生成 $I$ 的正则序列，参见《代数》引理
\ref{algebra-lemma-grothendieck-regular-sequence}。

\medskip\noindent
因此，若 $\Ker(S\to B)$ 也由正则序列生成，则
$\Ker(S'\to B)$ 同样如此，参见《代数详论》引理
\ref{more-algebra-lemma-join-koszul-regular-sequences} 和
\ref{more-algebra-lemma-noetherian-finite-all-equivalent}。

\medskip\noindent
反之，假设 $J=\Ker(S'\to B)$ 由正则序列生成。由于 $I$ 的生成元
$x_1,\ldots,x_c$ 在 $\mathfrak m_{S'}/\mathfrak m_{S'}^2$
中映为线性无关的元素，映射
$I/\mathfrak m_{S'}I\to J/\mathfrak m_{S'}J$ 是单射。
因此 $J$ 有一个极小生成元系
$x_1,\ldots,x_c,y_1,\ldots,y_d$。该序列是正则序列，从而
$y_1,\ldots,y_d$ 在 $S$ 中的像构成生成 $\Ker(S\to B)$ 的正则序列。
引理得证。
\end{proof}

\noindent
在下面的命题中，请注意：若 $B$ 在 $A$ 上具有有限 Tor 维数，
特别是若 $B$ 在 $A$ 上平坦，则关于 Tor 消失的条件成立。

\begin{proposition}
\label{dpa-proposition-avramov-map}
设 $A\to B$ 是诺特局部环之间的局部同态。下列条件等价：
\begin{enumerate}
\item $B$ 是完全交，且只有有限多个 $p$ 使
$\text{Tor}^A_p(B,A/\mathfrak m_A)$ 非零；
\item $A$ 是完全交，且 $A^\wedge\to B^\wedge$ 是上述意义下的
完全交同态。
\end{enumerate}
\end{proposition}

\begin{proof}
设 $F_\bullet\to A/\mathfrak m_A$ 是由有限自由 $A$-模组成的分辨。
注意 $\text{Tor}^A_p(B,A/\mathfrak m_A)$ 是复形
$F_\bullet\otimes_AB$ 的第 $p$ 个同调。令完备化
$F_\bullet^\wedge=F_\bullet\otimes_AA^\wedge$。则
$F_\bullet^\wedge$ 是由有限自由 $A^\wedge$-模组成的
$A^\wedge/\mathfrak m_{A^\wedge}$ 的分辨（因为 $A\to A^\wedge$
平坦且有限模上的完备化正合，参见《代数》引理
\ref{algebra-lemma-completion-tensor} 和
\ref{algebra-lemma-completion-flat}）。因此
$$
F_\bullet^\wedge \otimes_{A^\wedge} B^\wedge =
F_\bullet \otimes_A B \otimes_B B^\wedge
$$
由 $B\to B^\wedge$ 的平坦性，得到
$$
\text{Tor}^{A^\wedge}_p(B^\wedge, A^\wedge/\mathfrak m_{A^\wedge}) =
\text{Tor}^A_p(B, A/\mathfrak m_A) \otimes_B B^\wedge
$$
由此可见，局部环同态 $A\to B$ 的条件 (1) 等价于局部环同态
$A^\wedge\to B^\wedge$ 的相同条件。因此可以假设 $A$、$B$
是完备诺特局部环（其余条件本来就是用完备化表述的）。

\medskip\noindent
假设 $A$、$B$ 是完备诺特局部环。选取《代数详论》引理
\ref{more-algebra-lemma-embed-map-Noetherian-complete-local-rings}
中的交换图
$$
\xymatrix{
S \ar[r] & B \\
R \ar[u] \ar[r] & A \ar[u]
}
$$
令 $I=\Ker(R\to A)$、$J=\Ker(S\to B)$。
命题现由引理 \ref{dpa-lemma-perfect-map-ci} 得到。
\end{proof}

\begin{remark}
\label{dpa-remark-no-good-ci-map}
对于一般诺特环之间的同态，似乎很难定义一个良好的“局部完全交同态”
概念。原因是：对诺特局部环 $A$，$A\to A^\wedge$ 的纤维不是局部
完全交环。因此，若 $A\to B$ 是诺特局部环之间的局部同态，且完备化
之间的映射 $A^\wedge\to B^\wedge$ 是上述意义下的完全交同态，
那么一般而言，$(A_\mathfrak p)^\wedge\to(B_\mathfrak q)^\wedge$
{\bf 不是}上述意义下的完全交同态。只研究优秀诺特环可以解决这个问题。
更一般地，可以研究形式纤维为完全交的诺特环，参见
\cite{Rodicio-ci}。我们将在《对偶化复形》第
\ref{dualizing-section-formal-fibres} 节发展这一理论。
\end{remark}

\noindent
作为本节的结尾，我们把上述概念与《代数详论》第
\ref{dpa-section-lci} 节中引入的概念作比较。

\begin{lemma}
\label{dpa-lemma-well-defined-if-you-can-find-good-factorization}
考虑诺特局部环的交换图
$$
\xymatrix{
S \ar[r] & B \\
& A \ar[lu] \ar[u]
}
$$
，其中 $S\to B$ 为满射，$A\to S$ 平坦，且
$S/\mathfrak m_AS$ 是正则局部环。下列条件等价：
\begin{enumerate}
\item $\Ker(S\to B)$ 由正则序列生成；
\item $A^\wedge\to B^\wedge$ 是上述定义的完全交同态。
\end{enumerate}
\end{lemma}

\begin{proof}
从略。提示：证明与引理 \ref{dpa-lemma-ci-map-well-defined} 的证明第一段
中的论证相同。
\end{proof}

\begin{lemma}
\label{dpa-lemma-finite-type-lci-map}
设 $A$ 是诺特环，$A\to B$ 是有限型环同态。下列条件等价：
\begin{enumerate}
\item $A\to B$ 是《代数详论》定义
\ref{more-algebra-definition-local-complete-intersection} 意义下的
局部完全交；
\item 对每个素理想 $\mathfrak q\subset B$，令
$\mathfrak p=A\cap\mathfrak q$，则环同态
$(A_\mathfrak p)^\wedge\to(B_\mathfrak q)^\wedge$ 是上述意义下的
完全交同态。
\end{enumerate}
\end{lemma}

\begin{proof}
选取满射 $R=A[x_1,\ldots,x_n]\to B$。注意 $A\to R$ 平坦且纤维正则。
设 $I$ 是 $R\to B$ 的核。假设 (2)。由引理
\ref{dpa-lemma-well-defined-if-you-can-find-good-factorization} 和《代数》引理
\ref{algebra-lemma-regular-sequence-in-neighbourhood}，$I$ 局部由正则序列
生成。换言之，(1) 成立。反之，假设 (1)。将 $R$ 和 $B$ 局部化后，
可假设 $I$ 由 Koszul 正则序列生成。由《代数详论》引理
\ref{more-algebra-lemma-noetherian-finite-all-equivalent}，$I$ 局部由
正则序列生成。故由引理
\ref{dpa-lemma-well-defined-if-you-can-find-good-factorization}，(2) 成立。
略去若干细节。
\end{proof}

\begin{lemma}
\label{dpa-lemma-avramov-map-finite-type}
设 $A$ 是诺特环，$A\to B$ 是有限型环同态，并且
$\Spec(B)\to\Spec(A)$ 的像包含 $\Spec(A)$ 的所有闭点。
下列条件等价：
\begin{enumerate}
\item $B$ 是完全交，且 $A\to B$ 具有有限 Tor 维数；
\item $A$ 是完全交，且 $A\to B$ 是《代数详论》定义
\ref{more-algebra-definition-local-complete-intersection} 意义下的
局部完全交。
\end{enumerate}
\end{lemma}

\begin{proof}
这是命题 \ref{dpa-proposition-avramov-map} 经由引理
\ref{dpa-lemma-finite-type-lci-map} 得到的改述。细节从略。
\end{proof}








\section{光滑环同态与对角}
\label{dpa-section-smooth-diagonal-perfect}

\noindent
本节利用上述材料，将光滑环同态刻画为对角完美的环同态。

\begin{lemma}
\label{dpa-lemma-local-perfect-diagonal}
设 $A\to B$ 是诺特局部环之间的局部同态，并且 $B$ 在 $A$ 上平坦且
本质有限型。若
$$
B \otimes_A B \longrightarrow B
$$
是完美环同态，即 $B$ 在 $B\otimes_AB$ 上具有有限 Tor 维数，
则 $B$ 是某个光滑 $A$-代数的局部化。
\end{lemma}

\begin{proof}
由于 $B$ 在 $A$ 上本质有限型，$B\otimes_AB$ 也如此；特别地，
$B\otimes_AB$ 是诺特环。因此，$B\otimes_AB$ 的商 $B$ 在
$B\otimes_AB$ 上伪凝聚（《代数详论》引理
\ref{more-algebra-lemma-Noetherian-pseudo-coherent}），这说明环同态的
完美性（《代数详论》定义
\ref{more-algebra-definition-pseudo-coherent-perfect}）与有限 Tor 维数
条件一致。

\medskip\noindent
可将 $B$ 写成 $R/K$，其中 $R$ 是 $A[x_1,\ldots,x_n]$ 在某个素理想处的
局部化，$K\subset R$ 是理想。记 $\mathfrak m\subset R\otimes_AR$
为极大理想，它是 $B$ 的极大理想在满射
$R\otimes_AR\to B\otimes_AB\to B$ 下的逆像。于是有满射
$$
(R \otimes_A R)_\mathfrak m \to (B \otimes_A B)_\mathfrak m \to B
$$
，从而得到引理 \ref{dpa-lemma-injective} 中的理想
$I\subset J\subset(R\otimes_AR)_\mathfrak m$。因此
$I/\mathfrak mI\to J/\mathfrak mJ$ 是单射。

\medskip\noindent
将 $K$ 写成 $(f_1,\ldots,f_r)$，其中 $r$ 极小。可以并确实假设
$f_i\in R$ 是 $A[x_1,\ldots,x_n]$ 中某个也记作 $f_i$ 的元素之像。
注意 $I$ 由 $f_1\otimes1,\ldots,f_r\otimes1$ 和
$1\otimes f_1,\ldots,1\otimes f_r$ 生成。我们断言这是 $I$ 的一个
极小生成元集。事实上，若 $\kappa$ 是 $R$、$B$、
$(R\otimes_AR)_\mathfrak m$ 和 $(B\otimes_AB)_\mathfrak m$ 的共同剩余域，
则有映射
$R \otimes_A R \to R \otimes_A \kappa \oplus \kappa \otimes_A R$
，它经由 $(R\otimes_AR)_\mathfrak m$ 分解。由于 $B$ 在 $A$ 上平坦，
并且有短正合列 $0\to K\to R\to B\to0$，故
$K\otimes_A\kappa\subset R\otimes_A\kappa$，参见《代数》引理
\ref{algebra-lemma-flat-tor-zero}。因此，将映射
$(R \otimes_A R)_\mathfrak m \to R \otimes_A \kappa \oplus \kappa \otimes_A R$
限制到 $I$，得到映射
$$
I \to K \otimes_A \kappa \oplus \kappa \otimes_A K \to
K \otimes_B \kappa \oplus \kappa \otimes_B K.
$$
元素
$f_1 \otimes 1, \ldots, f_r \otimes 1, 1 \otimes f_1, \ldots, 1 \otimes f_r$
映为该映射目标的一组基，因为由 Nakayama 引理（《代数》引理
\ref{algebra-lemma-NAK}），$f_1,\ldots,f_r$ 映为
$K\otimes_B\kappa$ 的一组基。这证明了断言。

\medskip\noindent
理想 $J$ 由 $f_1\otimes1,\ldots,f_r\otimes1$ 以及元素
$x_1\otimes1-1\otimes x_1,\ldots,x_n\otimes1-1\otimes x_n$
生成（对证明而言，只须知道这些元素属于理想 $J$）。现在可写成
$$
f_i \otimes 1 - 1 \otimes f_i =
\sum g_{ij} (x_j \otimes 1 - 1 \otimes x_j)
$$
其中 $g_{ij}\in(R\otimes_AR)_\mathfrak m$。这是关于
$A[x_1,\ldots,x_n]$ 中元素的一般事实，证明从略。记
$a_{ij}\in\kappa$ 为 $g_{ij}$ 的像。另一项计算表明，$a_{ij}$ 是
$\partial f_i/\partial x_j$ 在 $\kappa$ 中的像。
$I/\mathfrak mI\to J/\mathfrak mJ$ 的单射性与上述评注迫使矩阵
$(a_{ij})$ 具有极大秩 $r$。令
$$
C = A[x_1, \ldots, x_n]/(f_1, \ldots, f_r)
$$
并考虑朴素余切复形
$\NL_{C/A}\cong(C^{\oplus r}\to C^{\oplus n})$，其中映射由偏导数矩阵
给出。因此，$\NL_{C/A}\otimes_AB$ 同构于放在次数 $0$ 的秩
$n-r$ 自由 $B$-模。故对某个在 $B$ 中映为单位的 $g\in C$，
$C_g$ 在 $A$ 上光滑，参见《代数》引理
\ref{algebra-lemma-smooth-at-point}。证明完成。
\end{proof}

\begin{lemma}
\label{dpa-lemma-perfect-diagonal}
设 $A\to B$ 是诺特环之间的平坦有限型环同态。若
$$
B \otimes_A B \longrightarrow B
$$
是完美环同态，即 $B$ 在 $B\otimes_AB$ 上具有有限 Tor 维数，
则 $B$ 是光滑 $A$-代数。
\end{lemma}

\begin{proof}
这由引理 \ref{dpa-lemma-local-perfect-diagonal} 以及光滑环同态的一般事实
得到，参见《代数》引理 \ref{algebra-lemma-smooth-at-point} 和
\ref{algebra-lemma-locally-smooth}。或者，读者可以略微修改引理
\ref{dpa-lemma-local-perfect-diagonal} 的证明来证明本引理。
\end{proof}






\section{余法模的自由性}
\label{dpa-section-freeness-conormal}

\noindent
Tate 分辨及其上的导子可用于证明本节结果的（更强）版本，
参见 \cite{Iyengar}。另外两篇更初等的参考文献是
\cite{Vasconcelos} 和 \cite{Ferrand-lci}。

\begin{lemma}
\label{dpa-lemma-free-summand-in-ideal-finite-proj-dim}
\begin{reference}
\cite{Vasconcelos}
\end{reference}
设 $R$ 是诺特局部环，$I\subset R$ 是在 $R$ 上具有有限投射维数的理想。
若 $F\subset I/I^2$ 是同构于 $R/I$ 的直和项，则存在非零因子
$x\in I$，使得 $x$ 在 $I/I^2$ 中的像生成 $F$。
\end{lemma}

\begin{proof}
由假设，可选取有限自由分辨
$$
0 \to R^{\oplus n_e} \to R^{\oplus n_{e-1}} \to \ldots \to
R^{\oplus n_1} \to R \to R/I \to 0
$$
于是 $\varphi_1:R^{\oplus n_1}\to R$ 的秩为 $1$，并由《代数》命题
\ref{algebra-proposition-what-exact} 可知 $I$ 包含一个非零因子 $y$。
设 $\mathfrak p_1,\ldots,\mathfrak p_n$ 是 $R$ 的伴随素理想，
参见《代数》引理 \ref{algebra-lemma-finite-ass}。选取理想
$I^2\subset J\subset I$，使得 $J/I^2=F$。由于 $y^2\in J$ 且
$y^2\notin\mathfrak p_i$，对所有 $i$ 有
$J\not\subset\mathfrak p_i$，参见《代数》引理
\ref{algebra-lemma-ass-zero-divisors}。由 Nakayama 引理（《代数》引理
\ref{algebra-lemma-NAK}），$J\not\subset\mathfrak mJ+I^2$。
由《代数》引理 \ref{algebra-lemma-silly}，可选取 $x\in J$，使得
$x\notin\mathfrak mJ+I^2$，且对 $i=1,\ldots,n$，
$x\notin\mathfrak p_i$。于是 $x$ 是非零因子，并且 $x$ 在 $I/I^2$
中的像由 Nakayama 引理生成直和项 $J/I^2\cong R/I$。
\end{proof}

\begin{lemma}
\label{dpa-lemma-vasconcelos}
\begin{reference}
\cite[定理 1.1]{Vasconcelos} 的局部版本
\end{reference}
设 $R$ 是诺特局部环，$I\subset R$ 是在 $R$ 上具有有限投射维数的理想。
若 $F\subset I/I^2$ 是秩为 $r$ 的自由直和项，则存在正则序列
$x_1,\ldots,x_r\in I$，使得
$x_1\bmod I^2,\ldots,x_r\bmod I^2$ 生成 $F$。
\end{lemma}

\begin{proof}
若 $r=0$，则无须证明。假设 $r>0$。由引理
\ref{dpa-lemma-free-summand-in-ideal-finite-proj-dim}，可选取 $x\in I$，
使得 $x$ 是非零因子，且 $x\bmod I^2$ 生成 $F$ 中一个同构于
$R/I$ 的直和项。考虑环 $R'=R/(x)$ 和理想 $I'=I/(x)$。
显然 $R'/I'=R/I$。短正合列
$$
0 \to R/I \xrightarrow{x} I/xI \to I' \to 0
$$
分裂，因为映射 $I/xI\to I/I^2$ 将 $xR/xI$ 映为一个直和项。
现在 $I/xI=I\otimes_R^\mathbf{L}R'$ 在 $R'$ 上具有有限投射维数，
参见《代数详论》引理 \ref{more-algebra-lemma-perfect-module} 和
\ref{more-algebra-lemma-pull-perfect}。故直和项 $I'$ 在 $R'$ 上具有
有限投射维数。另一方面，有短正合列
$0\to xR/xI\to I/I^2\to I'/(I')^2\to0$，从而
$I'/(I')^2$ 具有秩为 $r-1$ 的自由直和项
$F'=F/(R/I\cdot x)$。对 $r$ 作归纳，可选取正则序列
$x'_2,\ldots,x'_r\in I'$，使其同余类自由生成 $F'$。
令 $x_1=x$，并任取 $R$ 中提升 $x'_1,\ldots,x'_r$ 的元素
$x_2,\ldots,x_r$，即可得到引理结论。
\end{proof}

\begin{proposition}
\label{dpa-proposition-regular-ideal}
\begin{reference}
\cite[推论 1]{Vasconcelos} 的一个变体。另见
\cite{Iyengar} 和 \cite{Ferrand-lci}。
\end{reference}
设 $R$ 是诺特环，$I\subset R$ 是具有有限投射维数的理想，并且
$I/I^2$ 在 $R/I$ 上有限局部自由。则 $I$ 是正则理想
（《代数详论》定义 \ref{more-algebra-definition-regular-ideal}）。
\end{proposition}

\begin{proof}
由《代数》引理 \ref{algebra-lemma-regular-sequence-in-neighbourhood}，
只须证明对每个 $\mathfrak p\supset I$，
$I_\mathfrak p\subset R_\mathfrak p$ 由正则序列生成。
故可假设 $R$ 是局部环。若 $I/I^2$ 的秩为 $r$，则由引理
\ref{dpa-lemma-vasconcelos}，存在生成 $I/I^2$ 的正则序列
$x_1,\ldots,x_r\in I$。由 Nakayama 引理（《代数》引理
\ref{algebra-lemma-NAK}），$I$ 由 $x_1,\ldots,x_r$ 生成。
\end{proof}

\noindent
对任意环的局部完全交同态 $A\to B$，朴素余切复形 $\NL_{B/A}$
是 Tor 振幅位于 $[-1,0]$ 的完美复形，参见《代数详论》引理
\ref{more-algebra-lemma-lci-NL}。利用上述结果，可以证明这在某些情形
刻画局部完全交同态。

\begin{lemma}
\label{dpa-lemma-perfect-NL-lci}
设 $A\to B$ 是诺特环之间的完美环同态（《代数详论》定义
\ref{more-algebra-definition-pseudo-coherent-perfect}）。下列条件等价：
\begin{enumerate}
\item $\NL_{B/A}$ 的 Tor 振幅位于 $[-1,0]$；
\item $\NL_{B/A}$ 是 $D(B)$ 中 Tor 振幅位于 $[-1,0]$ 的完美对象；
\item $A\to B$ 是局部完全交（《代数详论》定义
\ref{more-algebra-definition-local-complete-intersection}）。
\end{enumerate}
\end{lemma}

\begin{proof}
将 $B$ 写成 $A[x_1,\ldots,x_n]/I$。则 $\NL_{B/A}$ 由 $B$-模复形
$$
I/I^2 \longrightarrow \bigoplus B \text{d}x_i
$$
表示，其中 $I/I^2$ 放在次数 $-1$。由于次数 $0$ 项有限自由，
该复形的 Tor 振幅位于 $[-1,0]$，当且仅当 $I/I^2$ 是平坦 $B$-模，
参见《代数详论》引理 \ref{more-algebra-lemma-last-one-flat}。
由于 $I/I^2$ 是有限 $B$-模且 $B$ 诺特，这又等价于 $I/I^2$
是有限局部自由 $B$-模（《代数》引理
\ref{algebra-lemma-finite-projective}）。因此 (1) 与 (2) 的等价性显然。
此外，应用命题 \ref{dpa-proposition-regular-ideal}，并注意正则理想既是
Koszul 正则理想也是拟正则理想（参见《代数详论》第
\ref{more-algebra-section-ideals} 节），也得到 (1) 与 (3) 的等价性。
\end{proof}

\begin{lemma}
\label{dpa-lemma-flat-fp-NL-lci}
设 $A\to B$ 是有限表示的平坦环同态。下列条件等价：
\begin{enumerate}
\item $\NL_{B/A}$ 的 Tor 振幅位于 $[-1,0]$；
\item $\NL_{B/A}$ 是 $D(B)$ 中 Tor 振幅位于 $[-1,0]$ 的完美对象；
\item $A\to B$ 是 syntomic 同态（《代数》定义
\ref{algebra-definition-lci}）；
\item $A\to B$ 是局部完全交（《代数详论》定义
\ref{more-algebra-definition-local-complete-intersection}）。
\end{enumerate}
\end{lemma}

\begin{proof}
(3) 与 (4) 的等价性是《代数详论》引理
\ref{more-algebra-lemma-syntomic-lci}。

\medskip\noindent
若 $A\to B$ 是 syntomic 同态，则可找到余笛卡尔图
$$
\xymatrix{
B_0 \ar[r] & B \\
A_0 \ar[r] \ar[u] & A \ar[u]
}
$$
，使得 $A_0\to B_0$ 是 syntomic 同态且 $A_0$ 诺特，参见《代数》引理
\ref{algebra-lemma-limit-module-finite-presentation} 和
\ref{algebra-lemma-colimit-lci}。由引理 \ref{dpa-lemma-perfect-NL-lci}，
$\NL_{B_0/A_0}$ 是 Tor 振幅位于 $[-1,0]$ 的完美对象。
由《代数详论》引理 \ref{more-algebra-lemma-base-change-NL-flat}，
$\NL_{B/A}=\NL_{B_0/A_0}\otimes_{B_0}^\mathbf{L}B$ 也具有同样性质
（另见《代数详论》引理 \ref{more-algebra-lemma-pull-tor-amplitude} 和
\ref{more-algebra-lemma-pull-perfect}）。这证明 (3) 蕴含 (2)。

\medskip\noindent
假设 (1)。由《代数详论》引理
\ref{more-algebra-lemma-base-change-NL-flat}，对每个以域 $k$ 为目标的
环同态 $A\to k$，$\NL_{B\otimes_Ak/k}$ 的 Tor 振幅位于
$[-1,0]$（参见《代数详论》引理
\ref{more-algebra-lemma-pull-tor-amplitude}）。故由引理
\ref{dpa-lemma-perfect-NL-lci}，$k\to B\otimes_Ak$ 是局部完全交同态。
于是按定义 $A\to B$ 是 syntomic 同态。这证明 (1) 蕴含 (3)，
并完成证明。
\end{proof}





\section{Koszul 复形与 Tate 分辨}
\label{dpa-section-koszul-vs-tate}

\noindent
本节把《代数详论》引理
\ref{more-algebra-lemma-sequence-Koszul-complexes} 的结论“提升”到
带有与微分分次结构相容的除幂的微分分次代数范畴中（请注意，本节把
Koszul 复形表示为链复形，而上述引文中使用的是上链复形）。

\medskip\noindent
设 $R$ 是环。设 $I \subset R$ 是由 $f_1, \ldots, f_r \in R$ 生成的
理想。对 $n \geq 1$，记
$$
K_n = K_{n, \bullet} = R\langle \xi_1, \ldots, \xi_r\rangle
$$
为微分分次 Koszul 代数，其中 $\xi_i$ 的次数为 $1$，且
$\text{d}(\xi_i) = f_i^n$。其上存在唯一的除幂结构（如定义
\ref{dpa-definition-divided-powers-dga} 中所述），见例
\ref{dpa-example-adjoining-odd}。对 $m > n$，转移映射
$K_m \to K_n$ 是与除幂相容的微分分次代数映射，它把 $\xi_i$
送到 $f_i^{m - n}\xi_i$。

\begin{lemma}
\label{dpa-lemma-lift-tate-to-koszul}
在上述情形中，若 $R$ 是诺特环，则对每个 $n$，存在 $N \geq n$
以及映射
$$
K_N \to A \to R/(f_1^N, \ldots, f_r^N)\quad\text{以及}\quad A \to K_n
$$
满足下列性质：
\begin{enumerate}
\item $(A, \text{d}, \gamma)$ 如定义
\ref{dpa-definition-divided-powers-dga} 所述；
\item $A \to R/(f_1^N, \ldots, f_r^N)$ 是拟同构；
\item 复合 $K_N \to A \to R/(f_1^N, \ldots, f_r^N)$ 是典范映射；
\item 复合 $K_N \to A \to K_n$ 是转移映射；
\item $A_0 = R \to R/(f_1^N, \ldots, f_r^N)$ 是典范满射；
\item $A$ 是 $R$ 上的分次除幂多项式代数，且每个次数中只有有限多个
生成元；以及
\item $A \to K_n$ 是与除幂相容的微分分次 $R$-代数同态，并在
$0$ 次同调上诱导典范映射
$R/(f_1^N, \ldots, f_r^N) \to R/(f_1^n, \ldots, f_r^n)$。
\end{enumerate}
条件 (4) 意味着，$A$ 是从 $A_0$ 出发，依照例
\ref{dpa-example-adjoining-odd} 或 \ref{dpa-example-adjoining-even}，
在每个次数 $> 0$ 中逐次添加有限集变量 $T$ 构造而成的。
\end{lemma}

\begin{proof}
固定 $n$。若 $r = 0$，则可直接取 $A = R$。假设 $r > 0$。由
《代数详论》引理 \ref{more-algebra-lemma-sequence-Koszul-complexes}
（将其改写为链复形的语言），可选取
$$
n_{r} > n_{r - 1} > \ldots > n_1 > n_0 = n
$$
使得 Koszul 代数上的转移映射 $K_{n_{i + 1}} \to K_{n_i}$（见上）
在次数 $> 0$ 的同调上诱导零映射。我们将取 $N = n_r$ 来证明本引理。

\medskip\noindent
我们完全按照引理 \ref{dpa-lemma-tate-resolution} 的陈述与证明来构造
$A$。于是将有
$$
A = \colim A(m),\quad\text{且}\quad
A(0) \to A(1) \to A(2) \to \ldots \to R/(f_1^N, \ldots, f_r^N)
$$
这立即给出性质 (1)、(2)、(5) 和 (6)。为完成证明，我们将构造
$R$-代数映射 $K_N \to A \to K_n$。具体构造如下：
\begin{enumerate}
\item 同构 $A(1) \to K_N = K_{n_r}$；
\item 映射 $A(2) \to K_{n_{r - 1}}$；
\item $\ldots$
\item 映射 $A(r) \to K_{n_1}$；
\item 映射 $A(r + 1) \to K_{n_0} = K_n$；以及
\item 映射 $A \to K_n$。
\end{enumerate}
在每一步中，所构造的映射都位于以相容方式赋予除幂的微分分次代数之间；
每个映射都经由 Koszul 代数之间的转移映射与前一个映射相容，并在
$0$ 次同调上诱导显然的典范映射。

\medskip\noindent
回忆 $A(0) = R$。当 $m = 1$ 时，引理
\ref{dpa-lemma-tate-resolution} 的证明选取
$A(1) = R\langle T_1, \ldots, T_r\rangle$，其中 $T_i$ 的次数为
$1$，且 $\text{d}(T_i) = f_i^N$。也就是说，$f_i^N$ 是
$A(0) \to R/(f_1^N, \ldots, f_r^N)$ 的核的生成元。因此，对
$A(1) \to K_N = K_{n_r}$ 使用映射
$$
\varphi_1 : A(1) \longrightarrow K_{n_r},\quad T_i \longmapsto \xi_i
$$
它是同构。

\medskip\noindent
当 $m = 2$ 时，引理 \ref{dpa-lemma-tate-resolution} 的证明中的构造选取
生成元 $e_1, \ldots, e_t \in \Ker(\text{d} : A(1)_1 \to A(1)_0)$。
接着取
$A(2) = A(1)\langle T_1, \ldots, T_t\rangle$
作为除幂多项式代数，其中 $T_i$ 的次数为 $2$，且
$\text{d}(T_i) = e_i$。由于 $\varphi_1(e_i)$ 是 $K_{n_r}$ 中的
闭链，由上述对 $n_r$ 与 $n_{r - 1}$ 的选取可知，它在
$K_{n_{r - 1}}$ 中的像是边界。因此可以构造交换图
$$
\xymatrix{
A(1) \ar[d] \ar[r]_{\varphi_1} & K_{n_r} \ar[d] \\
A(2) \ar[r]^{\varphi_2} & K_{n_{r - 1}}
}
$$
方法是把 $T_i$ 送到次数 $2$ 中边界等于 $\varphi_1(e_i)$ 之像的
一个元素。由除幂多项式代数的泛性质，映射 $\varphi_2$ 存在，并与
微分和除幂相容。

\medskip\noindent
对 $m = 2, \ldots, r$，完全以相同方式构造代数 $A(m)$ 与映射
$\varphi_m : A(m) \to K_{n_{r + 1 - m}}$。

\medskip\noindent
给定映射 $A(r) \to K_{n_1}$，可见复合
$H_r(A(r)) \to H_r(K_{n_1}) \to H_r(K_{n_0}) \subset (K_{n_0})_r$
为零，故可将其延拓为 $A(r + 1) \to K_{n_0} = K_n$，只需把
$A(r + 1)$ 的新多项式生成元送到零。

\medskip\noindent
构造出 $A(r + 1) \to K_{n_0} = K_n$ 后，只需把新的多项式生成元
送到零，即可按唯一可能的方式延拓到
$A(r + 2), A(r + 3), \ldots$。证明完毕。
\end{proof}

\begin{remark}
\label{dpa-remark-pro-system-koszul}
在上述情形中，若 $R$ 是诺特环，则可归纳地选取整数列
$1 = n_0 < n_1 < n_2 < \ldots $，使得对
$i = 1, 2, 3, \ldots$，存在引理
\ref{dpa-lemma-lift-tate-to-koszul} 中的映射
$K_{n_i} \to A_i \to R/(f_1^{n_i}, \ldots, f_r^{n_i})$
与 $A_i \to K_{n_{i - 1}}$。把复合
$A_{i + 1} \to K_{n_i} \to A_i$ 记作 $A_{i + 1} \to A_i$。
则图
$$
\xymatrix{
K_{n_1} \ar[d] &
K_{n_2} \ar[d] \ar[l] &
K_{n_3} \ar[d] \ar[l] &
\ldots \ar[l] \\
A_1 \ar[d] &
A_2 \ar[l] \ar[d] &
A_3 \ar[l] \ar[d] &
\ldots \ar[l] \\
K_1 &
K_{n_1} \ar[l] &
K_{n_2} \ar[l] &
\ldots \ar[l]
}
$$
交换。由此可见，在带有相容除幂的微分分次 $R$-代数范畴中，逆系统
$(K_n)$ 与 $(A_n)$ 是 pro-同构的。
\end{remark}

\begin{lemma}
\label{dpa-lemma-compute-cohomology-adjoin-variable}
设 $(A, \text{d}, \gamma)$、$d \geq 1$、$f \in A_{d - 1}$ 以及
$A\langle T \rangle$ 如引理 \ref{dpa-lemma-extend-differential} 中所述。
\begin{enumerate}
\item 若 $d = 1$，则有长正合列
$$
\ldots \to H_0(A) \xrightarrow{f} H_0(A) \to H_0(A\langle T \rangle) \to 0
$$
\item 当 $d = 2$ 时，存在有界谱序列
$(E_1)_{i, j} = H_{j - i}(A) \cdot T^{[i]}$
收敛到 $H_{i + j}(A\langle T \rangle)$。微分
$(d_1)_{i, j} : H_{j - i}(A) \cdot T^{[i]} \to
H_{j - i + 1}(A) \cdot T^{[i - 1]}$
把 $\xi \cdot T^{[i]}$ 送到 $f \xi \cdot T^{[i - 1]}$ 的类。
\item 根据需要在此补充其他次数的情形。
\end{enumerate}
\end{lemma}

\begin{proof}
当 $d = 1$ 时，有复形的短正合列
$$
0 \to A \to A\langle T \rangle \to A \cdot T \to 0
$$
由此容易得到结论 (1)。当 $d = 2$ 时，把 $A\langle T \rangle$
视为带有子复形
$$
F^pA\langle T \rangle = \bigoplus\nolimits_{i \leq p} A \cdot T^{[i]}
$$
的滤过链复形。应用《同调代数》第
\ref{homology-section-filtered-complex} 节的谱序列（将其改写为链复形）
即得 (2)。
\end{proof}

\noindent
稍后将用到下述引理。

\begin{lemma}
\label{dpa-lemma-construct-some-maps}
在上述情形中，对所有 $n \geq t \geq 1$，存在 $N > n$ 以及映射
$$
K_t \longrightarrow K_n \otimes_R K_t
$$
位于左微分分次 $K_N$-模的导出范畴中，使它与乘法映射的复合
（无论按哪个方向复合）都是转移映射。
\end{lemma}

\begin{proof}
先对 $r = 1$ 证明。令 $f = f_1$。写
$K_t = R\langle x \rangle$、
$K_n = R\langle y \rangle$ 与 $K_N = R\langle z \rangle$，其中
$x$、$y$、$z$ 的次数均为 $1$，且
$\text{d}(x) = f^t$、$\text{d}(y) = f^n$、$\text{d}(z) = f^N$。
我们断言，对所有 $N > t$，存在拟同构
$$
B_{N, t} = R\langle x, z, u \rangle
\longrightarrow
K_t,\quad
x \mapsto x,\quad
z \mapsto f^{N - t}x,\quad
u \mapsto 0
$$
这里左端表示变量 $x,z$ 的次数为 $1$、变量 $u$ 的次数为 $2$ 的
除幂多项式代数，并且
$\text{d}(x) = f^t$、$\text{d}(z) = f^N$、且
$\text{d}(u) = z - f^{N - t}x$。为证明该断言，注意
$H_*(R\langle x, z\rangle)$ 的下列三个子模相同：
\begin{enumerate}
\item $H_*(R\langle x, z\rangle) \to H_*(K_t)$ 的核；
\item 映射
$z - f^{N - t}x : H_*(R\langle x, z\rangle) \to H_*(R\langle x, z\rangle)$
的像；以及
\item 映射
$z - f^{N - t}x : H_*(R\langle x, z\rangle) \to H_*(R\langle x, z\rangle)$
的核。
\end{enumerate}
这一点可由直接计算证明\footnote{提示：令
$z' = z - f^{N - t}x$，则
$R\langle x, z\rangle = R\langle x, z'\rangle$ 且
$\text{d}(z') = 0$；此外，映射
$R\langle x, z'\rangle \to K_t$ 正是把 $z'$ 杀掉的映射。}，
计算从略。再应用引理
\ref{dpa-lemma-compute-cohomology-adjoin-variable} 的第 (2) 部分，
即可得出该断言。

\medskip\noindent
通过把 $z$ 送到 $z$ 的微分分次 $R$-代数同态
$K_N \to B_{N, t}$，可把 $B_{N, t} \to K_t$ 看作左微分分次
$K_N$-模的拟同构。为定义引理陈述中的箭头，我们使用同态
$$
B_{N, t} = R\langle x, z, u \rangle \to K_n \otimes_R K_t,\quad
x \mapsto 1 \otimes x,\quad
z \mapsto f^{N - n}y \otimes 1,\quad
u \mapsto - f^{N - n - t}y \otimes x
$$
只要假设 $N \geq n + t$，这就有意义。一个简单明快的计算表明，
它与乘法映射的（前复合或后复合）是转移映射。

\medskip\noindent
当 $r > 1$ 时，把每个 Koszul 代数写成单变量 Koszul 代数的张量积，
并应用上述构造。换言之，写
$$
K_t = R\langle x_1, \ldots, x_r\rangle =
R\langle x_1\rangle \otimes_R \ldots \otimes_R R\langle x_r\rangle
$$
其中 $x_i$ 的次数为 $1$，且 $\text{d}(x_i) = f_i^t$。在
$r > 1$ 的情形下，使用
$$
B_{N, t} =
R\langle x_1, z_1, u_1 \rangle
\otimes_R \ldots \otimes_R
R\langle x_r, z_r, u_r \rangle
$$
其中 $x_i,z_i$ 的次数为 $1$，$u_i$ 的次数为 $2$，并且
$\text{d}(x_i) = f_i^t$、$\text{d}(z_i) = f_i^N$、且
$\text{d}(u_i) = z_i - f_i^{N - t}x_i$。张量积映射
$B_{N, t} \to K_t$ 是拟同构，因为它是自由 $R$-模的上有界复形之间
诸拟同构的张量积。最后，将映射
$$
B_{N, t}
\to
K_n \otimes_R K_t =
R\langle y_1, \ldots, y_r\rangle \otimes_R
R\langle x_1, \ldots, x_r\rangle
$$
定义为 $r = 1$ 情形中所构造映射的张量积；也可直接用规则
$x_i \mapsto 1 \otimes x_i$、
$z_i \mapsto f_i^{N - n}y_i \otimes 1$ 与
$u_i \mapsto - f_i^{N - n - t}y_i \otimes x_i$ 来定义，只要
$N \geq n + t$，这些规则就有意义。细节从略。
\end{proof}
